\documentclass[11pt,twoside]{article}

\usepackage[margin=0.92in]{geometry}
\usepackage[T1]{fontenc}
\usepackage{lmodern}
\usepackage{microtype}
\usepackage{amsmath,amssymb,mathtools,amsthm}
\usepackage{enumitem}
\usepackage{xcolor}
\usepackage{hyperref}
\usepackage{fancyhdr}
\usepackage{titlesec}
\usepackage{longtable}
\usepackage{booktabs}
\usepackage{array}
\usepackage{tabularx}

\definecolor{FactorBlue}{RGB}{29,72,111}
\definecolor{FactorDark}{RGB}{38,43,48}
\definecolor{FactorGray}{RGB}{246,248,250}

\hypersetup{%
  colorlinks=true,
  linkcolor=FactorBlue,
  urlcolor=FactorBlue,
  pdftitle={Factor: A Construction, Proof, and Retained-Replay Kernel},
  pdfauthor={Foundation proposal},
  pdfsubject={A self-contained foundation proposal for Factor}
}

\fancypagestyle{factor}{%
  \fancyhf{}%
  \fancyhead[LO,LE]{Factor: A Construction, Proof, and Retained-Replay Kernel}%
  \fancyhead[RO,RE]{Proposal 6.1}%
  \fancyfoot[CO,CE]{\thepage}%
}
\fancypagestyle{plain}{%
  \fancyhf{}%
  \fancyhead[LO,LE]{Factor: A Construction, Proof, and Retained-Replay Kernel}%
  \fancyhead[RO,RE]{Proposal 6.1}%
  \fancyfoot[CO,CE]{\thepage}%
}
\pagestyle{factor}
\setlength{\headheight}{14pt}
\raggedbottom{}
\setlength{\emergencystretch}{3em}

\titleformat{\section}{\Large\bfseries\color{FactorDark}}{\thesection}{0.7em}{}
\titleformat{\subsection}{\large\bfseries\color{FactorDark}}{\thesubsection}{0.7em}{}
\titleformat{\subsubsection}{\normalsize\bfseries\color{FactorDark}}{\thesubsubsection}{0.7em}{}

\newcommand{\boxtext}[1]{%
\begin{center}
\fcolorbox{FactorBlue}{FactorGray}{%
\begin{minipage}{0.90\textwidth}
#1
\end{minipage}}
\end{center}}

\newcommand{\Apply}{\mathsf{Apply}}
\newcommand{\Closure}{\mathsf{Closure}}
\newcommand{\Reach}{\mathsf{Reach}}
\newcommand{\SameStructure}{\mathsf{SameStructure}}
\newcommand{\BoundaryCorresponds}{\mathsf{BoundaryCorresponds}}
\newcommand{\Derivable}{\mathsf{Derivable}}
\newcommand{\Novel}{\mathsf{Novel}}
\newcommand{\Family}{\mathsf{Family}}
\newcommand{\Necessary}{\mathsf{Necessary}}
\newcommand{\Formed}{\mathsf{Formed}}
\newcommand{\Replay}{\mathsf{Replay}}
\newcommand{\RealizedAt}{\mathsf{RealizedAt}}
\newcommand{\Proves}{\mathsf{Proves}}
\newcommand{\Derives}{\mathsf{Derives}}
\newcommand{\ValidDerivation}{\mathsf{ValidDerivation}}
\newcommand{\CandidateAt}{\mathsf{CandidateAt}}
\newcommand{\MatchAt}{\mathsf{MatchAt}}
\newcommand{\CompleteMatch}{\mathsf{CompleteMatch}}
\newcommand{\FiniteMember}{\mathsf{FiniteMember}}
\newcommand{\UniversalAdmissible}{\mathsf{UniversalAdmissible}}
\newcommand{\ProjectsAt}{\mathsf{ProjectsAt}}
\newcommand{\Rec}{\mathsf{Rec}}
\newcommand{\Fam}{\mathsf{Fam}}
\newcommand{\RefPos}{\mathsf{Ref}}
\newcommand{\ExactMatch}{\mathsf{ExactMatch}}

\newtheorem{theorem}{Theorem}
\newtheorem{lemma}[theorem]{Lemma}
\newtheorem{corollary}[theorem]{Corollary}

\title{\textbf{Factor: A Construction, Proof, and Retained-Replay Kernel}\\[0.35em]
\large A self-contained proposal}
\author{Foundation proposal}
\date{Proposal 6.1 --- August 15, 2026}

\begin{document}
\maketitle
\thispagestyle{empty}

\begin{abstract}
Factor gives structural meaning to finite constructions without granting force to
names, encodings, proof layouts, or storage choices.  Its normative input is an
abstract finite table over one closed structural signature; decoding bytes or host
objects into that table is not a Factor judgment.  Canonical encoding profiles
remain external maps but must satisfy stated totality, injectivity, round-trip, and
same-input obligations.  Every semantic relation is interpreted in one explicitly
fixed finite Factor model.  A closed catalogue
checks supplied finite derivation graphs independently of storage.  A separate
structural realization maps those derivations to retained fragments, and one replay
relation combines derivation validity, realization, and ancestry.

Semantic truth may be undecidable.  Reference checks nevertheless terminate
because they inspect only finite supplied structure.  Optional authority can select
exact claims and adopt external interpretations, but raw adoption neither proves a
claim nor chooses retained evidence.  Caches, indexes, creation procedures, and
physical storage remain nonsemantic.
\end{abstract}

\setcounter{tocdepth}{2}
\newpage
\tableofcontents
\newpage

\section*{Purpose and dependency boundary}
\addcontentsline{toc}{section}{Purpose and dependency boundary}

Factor answers:

\begin{quote}
Which earlier constructions produce this occurrence, and which inputs entered
without a retained derivation?
\end{quote}

An underived cause says only that no derivation was retained for that admission.
It does not assert truth, independence, novelty, authority, or external meaning.
An encoding or structural correspondence creates no commitment.  Only explicit
authority structure adopts or publishes an external interpretation.

Here \emph{encoding} means a byte string, database row, host object, or other
carrier outside the formal model.  The proposal necessarily assigns internal
Factor syntax a fixed meaning through the closed primitive signature and formulas
below.  Recognizing that formal syntax forms a Factor cause or judgment is not an
adoption of an external interpretation.

The proposal has four layers:

\begin{enumerate}[leftmargin=2.2em]
  \item semantic Factor relations;
  \item finite storage-independent derivations, conditional proof validity, and
        global proof validity;
  \item structural realization, retained fragments and selector links, and
        normative replay at a site; and
  \item creation, deployment, caching, indexing, and implementation work.
\end{enumerate}

\boxtext{%
\[
\begin{array}{c}
\text{semantic relations}\\[-0.1em]
\downarrow\\[-0.1em]
\text{finite structural derivations}\\[-0.1em]
\downarrow\\[-0.1em]
\text{retained realization and replay}\\[-0.1em]
\downarrow\\[-0.1em]
\text{implementation}
\end{array}
\]
No arrow points in the opposite direction.}

Meaning is fixed before evidence.  Proof can establish an already meaningful
judgment but cannot create its meaning.  Realization can retain a derivation but
cannot make that derivation valid.  Implementation can accelerate any fixed check
but cannot change its answer.

\section*{Metalevel: finite Factor models}
\addcontentsline{toc}{section}{Metalevel: finite Factor models}

\subsection*{One structural universe and its irreducible signature}

All positions and values in a normative Factor input inhabit one universe \(U\).
Occurrences, frontiers, contexts, frames, proof nodes, witnesses, authority
instances, and bridges are predicates or views selected by structure.  They are
not disjoint sorts.  A place accepts any value whose supplied structure satisfies
its fixed incident equations.

A normative finite input is the mathematical structure
\[
 G=(U_G,\operatorname{Own}_G,\operatorname{End}_G,
       \operatorname{Next}_G,\operatorname{Bind}_G),
 \qquad U_G\subseteq U,
\]
whose five components are literal finite tables, not procedures.  Equality is
ordinary equality in \(U\).  Call the resulting closed signature \(\Sigma_F\).
Its only primitive readings are:

\begin{itemize}[leftmargin=2em]
  \item \(\operatorname{Own}_G(a,p)\): \(p\) is an immediate position owned by
        \(a\);
  \item \(\operatorname{End}_G(p,x)\): incidence position \(p\) has endpoint
        \(x\);
  \item \(\operatorname{Next}_G(p,q)\): \(q\) is the explicitly ordered next
        sibling of \(p\); and
  \item \(\operatorname{Bind}_G(k,p)\): key value \(k\) is bound to exact
        graph position \(p\).
\end{itemize}

Every graph-backed owner, local position, value, and endpoint is in \(U_G\); thus
\(\operatorname{Pos}(G)=U_G\) when a carrier rather than a role is required.
Distinct incidence positions represent multiplicity even when their endpoints are
equal.  Direction and encoded order are represented by endpoint placement and
\(\operatorname{Next}_G\), not by labels.  Immediate ownership is functional in
its second component and acyclic.  An endpoint place has one endpoint, and
\(\operatorname{Next}_G\) relates only positions with the same immediate owner,
with at most one predecessor and successor.  The binding table is partial in its
first component.  A value is a key only because it occurs at a key place fixed by
one of the structural equations below.  Absence of a pair in an explicitly
restricted finite binding table is a failed lookup extensionally.

Supplied hypothetical structures use the same four structural relations
\(\operatorname{Own},\operatorname{End},\operatorname{Next},\operatorname{Bind}\)
on their finite positions; their graph-facing endpoints and binding results are
exact positions and tuples of \(G\).  They do not introduce a second syntax.

\subsection*{Exact record notation and the recognition boundary}

The following notation is normative.  It reduces every later record and pattern
claim to the primitive tables.  For distinct \(x_1,\ldots,x_n\),
\[
\begin{split}
 \Rec_G(r;x_1,\ldots,x_n)\quad\Longleftrightarrow\quad{}
 &\forall x\,[\operatorname{Own}_G(r,x)
             \leftrightarrow \textstyle\bigvee_{i=1}^{n}x=x_i]\\
 &{}\land
 \forall x,y\,[\operatorname{Own}_G(r,x)\land
                \operatorname{Own}_G(r,y)\Rightarrow\\[-0.4em]
 &\hspace{7.3em}(\operatorname{Next}_G(x,y)
      \leftrightarrow
      \textstyle\bigvee_{i=1}^{n-1}(x=x_i\land y=x_{i+1}))].
\end{split}
\]
For an exact finite set \(X\),
\[
 \Fam_G(r;X)\Longleftrightarrow
 \bigl(\forall x\,[\operatorname{Own}_G(r,x)\leftrightarrow x\in X]\bigr)
 \land\neg\exists x,y\in X\ \operatorname{Next}_G(x,y).
\]
An ordered family is a \(\Rec\) with its displayed finite sequence.  An optional
field is a \(\Fam\) of cardinality zero or one.  Finally,
\[
 \RefPos_G(q;x)\Longleftrightarrow
 \bigl(\forall y\,[\operatorname{End}_G(q,y)\leftrightarrow y=x]\bigr)
 \land\neg\exists s\,\operatorname{Own}_G(q,s).
\]
A pair record is fixed by
\[
 \operatorname{PairAt}_G(q;a,b)\Longleftrightarrow
 \exists k,v\,[\Rec_G(q;k,v)\land\RefPos_G(k;a)
                         \land\RefPos_G(v;b)].
\]
A finite relation or map is then fixed by
\[
\begin{split}
 \operatorname{RelationRecord}_G(r)\Longleftrightarrow{}&
 \exists X\,[\Fam_G(r;X)
   \land\forall q\in X\ \exists!a,b\,\operatorname{PairAt}_G(q;a,b)\\[-0.3em]
 &\quad{}\land\forall q,q'\in X\ \forall a,b\,[
   \operatorname{PairAt}_G(q;a,b)\land
   \operatorname{PairAt}_G(q';a,b)\Rightarrow q=q']],\\
 \operatorname{MapRecord}_G(r)\Longleftrightarrow{}&
 \operatorname{RelationRecord}_G(r)\\[-0.3em]
 &{}\land\forall q,q',a,b,b'\,[
   \operatorname{Own}_G(r,q)\land\operatorname{Own}_G(r,q')\\[-0.3em]
 &\hspace{5em}{}\land\operatorname{PairAt}_G(q;a,b)
   \land\operatorname{PairAt}_G(q';a,b')\Rightarrow q=q'].
\end{split}
\]
Its projected endpoint relation is not discretionary:
\[
 \operatorname{Pairs}_G(r)=
 \{(a,b):\exists q\,[\operatorname{Own}_G(r,q)
                   \land\operatorname{PairAt}_G(q;a,b)]\}.
\]
Whenever a later relation/map argument is denoted by \(V,\sigma,\mu,\alpha,m\),
the root remains part of its exact boundary and the mathematical relation is this
unique \(\operatorname{Pairs}\) projection.  These equations distinguish a
missing member, duplicate position, repeated endpoint, and encoded order without
a kind tag.

For a finite pattern \(K\) with internal positions \(I_K\), external parameters
\(A_K\), and a proposed map \(\eta\), write
\(\ExactMatch_G(K,\eta)\) exactly when \(\eta\) is one-to-one on \(I_K\), fixes
the supplied members of \(A_K\), and, for each
\(R\in\{\operatorname{Own},\operatorname{End},\operatorname{Next},
\operatorname{Bind}\}\),
\[
 R_G(\eta\vec x)\quad\Longleftrightarrow\quad R_K(\vec x)
\]
for every tuple incident with an internal position, including every
boundary-crossing tuple.  The image contains the whole owned closure of the pattern
and no other owned position.  Finite family variables in \(K\) range over the
exact displayed \(\Fam\), so this is a finite relational check, not matching by a
name or by an omitted tail.  ``Contains exactly'', ``exposes'', ``structural
pattern'', and ``complete/no-extra'' below are abbreviations for these
reflection equations.

The top-level grammar used later is fixed here.  Each row is completed by the
payload equations in the cited section; there is no additional recognizer.

\begin{center}
\begin{tabularx}{0.98\textwidth}{@{}>{\raggedright\arraybackslash}p{0.19\textwidth}
                                      >{\raggedright\arraybackslash}p{0.37\textwidth}
                                      X@{}}
\toprule
Structure & Exact root spine & Payload roots\\
\midrule
five-part body & \(\Rec_G(B;R,P,Q,T,N)\) & exact reference, place,
requirement, target, and declaration families\\
frontier & \(\Rec_G(f;P,e)\) & parent-reference family \(P\) and optional
entry reference \(e\)\\
base occurrence & \(\Rec_G(o;\kappa),\ \Rec_G(\kappa;B)\) & one formed body
record \(B\)\\
derived occurrence &
\(\Rec_G(o;\kappa),\ \Rec_G(\kappa;s,V,\sigma,\mu,\alpha)\) & one source
reference and four exact finite relation/map records\\
structural view & \(\Rec_G(E;A,C)\) & interface \(A\) and closed clause graph
\(C\)\\
finite relation & \(\Rec_G(E;A,T,e)\) & interface \(A\), exact tuple family
\(T\), and structural comparison boundary \(e\)\\
semantic closure & \(\Rec_G(E;A,S,t)\) & interface, seed, and
\(\RefPos_G(t;\operatorname{Step})\)\\
finite domain & \(\Rec_G(D;M)\) & exact member-reference family \(M\)\\
admissibility schema & \(\Rec_G(S;B,E,H,C,T,D)\) & binders, equations,
hypotheses, conclusion, fixed targets, and guards\\
retained envelope & \(\Rec_G(p;t,C,F,Q,L)\) & target reference, conclusion
boundary, fragments, path graph, and selector links\\
\bottomrule
\end{tabularx}
\end{center}

The structural-view clause grammar and every other payload grammar are the closed
finite alternatives displayed in their authoritative sections, each expressed by
\(\Rec\), \(\Fam\), \(\RefPos\), and \(\ExactMatch\).  A pattern name in prose
does not occur in \(G\) and cannot select a row.  Two simultaneous rows, no row,
or two matches with different field projections make the proposed structure
malformed.

This signature is the proposal's irreducible mathematical base.  Factor does
\emph{not} map bytes, database records, host objects, or wire tags into it, and it
cannot certify that two such adapters intend the same primitive tables.  Doing so
would require another interpretation function.  Section~\ref{sec:implementation}
therefore states the minimum external contract for a canonical encoding profile:
total deterministic parsing with explicit rejection, total injective encoding on
a declared domain, exact round trips, and a literal same-input criterion.  No
particular profile or parser is supplied here, and profile conformance remains an
external claim rather than a Factor judgment.

\subsection*{Five-part bodies and pure application}

A carried body is the finite record
\[
  B=(R,P,Q,T,N),
\]
meaning the exact spine \(\Rec(B;R,P,Q,T,N)\) in the supplied structural
tables.  Each of \(R,P,Q,T,N\) is an exact \(\Fam\) (or an explicitly ordered
\(\Rec\) where order is part of the body), and every member has the complete
endpoint pattern required by its collection.  In primitive notation, define
\[
 \operatorname{Empty}_G(q)\Longleftrightarrow
 \neg\exists x\,[\operatorname{Own}_G(q,x)\lor\operatorname{End}_G(q,x)],
\]
and let \(\operatorname{OrderedRefs}_G(r;x_1,\ldots,x_n)\) mean
\(\exists e_1,\ldots,e_n\,[\Rec_G(r;e_1,\ldots,e_n)\land
  \bigwedge_i\RefPos_G(e_i;x_i)]\) for \(n\ge1\), and
\(\operatorname{Empty}_G(r)\) for \(n=0\).  The five member predicates are
\[
\begin{aligned}
 \operatorname{ReferenceItem}_G(r)
  &\Longleftrightarrow
    \exists n\ge1,\vec x\ \operatorname{OrderedRefs}_G(r;\vec x),\\
 \operatorname{OpenItem}_G(p)
  &\Longleftrightarrow \operatorname{Empty}_G(p),\\
 \operatorname{RequirementItem}_G(q)
  &\Longleftrightarrow
    \exists O,E,X,\vec x\,[\Rec_G(q;O,E)
       \land\Fam_G(O;X)\\[-0.4em]
  &\hspace{6em}{}\land
    \forall a\in X\ \exists y\,\RefPos_G(a;y)
    \\[-0.4em]
  &\hspace{6em}{}\land\operatorname{OrderedRefs}_G(E;\vec x)
    \land(X\ne\varnothing\lor |\vec x|>0)],\\
 \operatorname{TargetItem}_G(t)
  &\Longleftrightarrow
    \exists e,x\,[\Rec_G(t;e)\land\RefPos_G(e;x)],\\
 \operatorname{NewItem}_G(n)
  &\Longleftrightarrow
    \exists p\,[\Rec_G(n;p)\land\operatorname{Empty}_G(p)].
\end{aligned}
\]
The two requirement fields are respectively its exact unordered open-endpoint
family and its ordered resolved-endpoint record; their union is nonempty.  Open references may end
at hypothetical places; resolved references end at exact supplied occurrences or
graph positions.  Collection membership fixes which displayed predicate applies,
so the same endpoint value may occur in several roles without a value sort.

Writing \(X_C\) for the exact child set of collection root \(C\), define
\[
\begin{split}
 \operatorname{BodyRecord}_G(B)\Longleftrightarrow
 \exists R,P,Q,T,N,X_R,X_P,X_Q,X_T,X_N\,[{}&
 \Rec_G(B;R,P,Q,T,N)\\
 &{}\land\bigwedge_C\Fam_G(C;X_C)\\
 &{}\land\forall r\in X_R\ \operatorname{ReferenceItem}_G(r)\\
 &{}\land\forall p\in X_P\ \operatorname{OpenItem}_G(p)\\
 &{}\land\forall q\in X_Q\ \operatorname{RequirementItem}_G(q)\\
 &{}\land\forall t\in X_T\ \operatorname{TargetItem}_G(t)\\
 &{}\land\forall n\in X_N\ \operatorname{NewItem}_G(n)].
\end{split}
\]
where \(R\) contains outgoing references, \(P\) open places, \(Q\) required links,
\(T\) introduced targets, and \(N\) new-place declarations.  Membership in these
five collections is primitive positional structure.  Each item retains its owner,
local position, direction where present, endpoints, multiplicity, aliasing, and
sharing.  No item carries a role name.

For a source body \(B_{\mathrm{src}}\), a finite family \(U_0\) of explicitly
supplied bodies, a partial assignment \(\sigma\) of source places, an injective
match \(\mu\) from every resolved requirement to a distinct matching reference,
and a fresh-position pairing \(\alpha\), define the graph-independent partial
function
\[
  \Apply(B_{\mathrm{src}},U_0,\sigma,\mu,\alpha)=B_{\mathrm{out}}.
\]
The domain of \(\alpha\) is exactly the unassigned places, new-place declarations,
unresolved requirements, and introduced targets.  Its range is exactly the fresh
output positions.  The source position determines the output collection; \(\alpha\)
contains no role field.

The operation is defined exactly when:

\begin{enumerate}[leftmargin=2.2em]
  \item every owner and position cited by \(U_0,\sigma,\mu,\alpha\) exists in the
        supplied finite material;
  \item \(\sigma\) is a partial function on \(P\);
  \item after assignment, \(\mu\) is total on the resolved requirements,
        injective, and selects references with the exact direction and endpoints;
  \item \(\alpha\) has the exact domain and range just stated and is one-to-one;
        and
  \item the candidate output contains all and only the items given by the output
        equations that follow.
\end{enumerate}

The output equations are:

\begin{itemize}[leftmargin=2em]
  \item an unassigned source place remains one ordinary open place;
  \item a new-place declaration becomes one ordinary open place and no declaration
        remains;
  \item a resolved requirement is discharged after its matched reference passes;
  \item an unresolved requirement remains, with every resolved endpoint retained
        and every open endpoint redirected through \(\alpha\);
  \item a resolved introduced target becomes one distinct outgoing reference from
        the result to the resolved occurrence; and
  \item an unresolved introduced target remains and points to its redirected
        result place.
\end{itemize}

Source references remain owned by the source and are not copied.  There is no
other output item.

\(\Apply\) reads no model, frontier, admission, proof, authority, storage, or
external meaning.  It terminates on fixed finite inputs and is deterministic once
\(\alpha\) is fixed.  Two fresh-position choices differ only by a bijective
renaming of result positions.

\subsection*{Model-formation conditions}

A \emph{Factor model} is a finite graph satisfying all of the following metalevel
conditions.

\begin{enumerate}[leftmargin=2.2em]
  \item Every value occupying a binding-key place has at most one binding.
  \item Every cited owner, local position, endpoint, and binding endpoint exists,
        and every finite record satisfies its displayed ownership and incidence
        equations.
  \item Each frontier satisfies \(\Rec_G(f;P,e)\), where \(P\) is its exact
        \(\Fam\) of parent-frontier references and \(e\) is an optional
        occurrence-reference \(\Fam\).  An
        initial frontier has empty \(P\) and no entry.  A noninitial frontier has
        nonempty \(P\) and exactly one entry.
  \item The frontier-parent graph is finite and acyclic.  One occurrence is the
        entry of at most one frontier.
  \item Every occurrence entry has exactly one of the two structurally disjoint
        cause shapes in Section~\ref{sec:kernel}: a base cause carrying one exact
        body, or a derived cause carrying one source, one complete selected
        occurrence/position map, and the maps \(\sigma,\mu,\alpha\).  The first
        shape has no application positions; the second has all and only the
        application positions.
  \item An occurrence entered by a frontier is absent from every strict parent
        development.  Every source, selected input, matched-reference owner,
        external endpoint, and other support used by its cause belongs to the
        required parent development.
  \item A base cause introduces its exact body.  A derived cause satisfies the
        authoritative body equation of Section~\ref{sec:kernel} against bodies
        obtained from strict ancestors.  Every admission and carried record
        satisfies the complete/no-extra equations stated there.
\end{enumerate}

These conditions restrict the domain over which all Factor relations are
interpreted.  They are not a Factor semantic relation, not a judgment, and not a
conclusion that proof or replay can establish.  A storage system may construct or
incrementally check a candidate graph, but those acts add no semantic fact.

\boxtext{\textbf{Fixed-model convention.}  For the rest of the authoritative
sections, fix one Factor model \(G\).  Every
semantic, proof, realization, and replay relation is relative to this same \(G\).
The parameter identifies immutable structure, not a repository handle, storage
address, realization site, or validator state.}

\section{Semantic kernel}\label{sec:kernel}

\subsection{Frontiers, developments, and admission}

Write
\[
\begin{split}
 \operatorname{FrontierRecord}_G(f;P,e,X_P,X_e)
 \Longleftrightarrow{}&\Rec_G(f;P,e)\land\Fam_G(P;X_P)
                  \land\Fam_G(e;X_e)\\
 &{}\land |X_e|\le1
 \land\forall q\in X_P\ \exists g\,\RefPos_G(q;g)\\
 &{}\land\forall q\in X_e\ \exists o\,\RefPos_G(q;o).
\end{split}
\]
Model formation requires one such projection for every frontier root and requires
all referenced \(g,o\) to satisfy their respective frontier and occurrence
equations.  Let \(\operatorname{Parents}_G(f)\) be the endpoints of the references
in \(X_P\), and let \(\operatorname{Entry}_G(f)\) be absent or the one endpoint
from \(X_e\).  For this unique projection, let \(\Delta_G(f)\) be the occurrence
entries in the finite parent closure of \(f\), including
\(\operatorname{Entry}_G(f)\) when present, and write
\[
  \Delta_G(P)=\bigcup_{g\in P}\Delta_G(g)
  \qquad(P\text{ an exact frontier set}).
\]
The relation
\[
  \operatorname{AdmittedAt}_G(o,f)
\]
holds exactly when
\[
 \exists P,e,X_P,X_e,q\,[
 \operatorname{FrontierRecord}_G(f;P,e,X_P,X_e)
 \land q\in X_e\land\RefPos_G(q;o)].
\]
An occurrence does
not carry a frontier back-reference.

The cause of an admitted occurrence \(o\) has exactly one of these shapes:
\[
\begin{array}{ll}
\text{base:}    & \operatorname{BaseCause}_G(o;B),\\[0.2em]
\text{derived:} & \operatorname{DerivedCause}_G
                    (o;o_{\mathrm{src}},V,\sigma,\mu,\alpha).
\end{array}
\]
These names abbreviate explicit substrate formulas, not a stored kind or a fixed
decoder.  In full,
\[
\begin{split}
 \operatorname{BaseCause}_G(o;B)\Longleftrightarrow{}&
 \exists\kappa\,[\Rec_G(o;\kappa)\land\Rec_G(\kappa;B)
                 \land\operatorname{BodyRecord}_G(B)],\\
 \operatorname{DerivedCause}_G
   (o;o_{\mathrm{src}},V,\sigma,\mu,\alpha)
 &\Longleftrightarrow
 \exists\kappa,s\,[\Rec_G(o;\kappa)
       \land\Rec_G(\kappa;s,V,\sigma,\mu,\alpha)\\[-0.4em]
 &\hspace{4.5em}{}\land\RefPos_G(s;o_{\mathrm{src}})
       \land\operatorname{RelationRecord}_G(V)\\[-0.4em]
 &\hspace{4.5em}{}\land\operatorname{MapRecord}_G(\sigma)
       \land\operatorname{MapRecord}_G(\mu)
       \land\operatorname{MapRecord}_G(\alpha)].
\end{split}
\]
Here \(\operatorname{BodyRecord}_G(B)\) is exactly the five-field \(\Rec\) and
its five complete family equations above.  The predicates
\(\operatorname{RelationRecord}\) and
\(\operatorname{MapRecord}\) are exactly the \(\Fam\)-of-pairs equations from the
substrate, with map-key functionality for the latter.  The first formula has one
cause field and therefore no application position.  The second has all and only
the displayed source, selected relation, assignment, reference match, and position
pairing fields; it has no copied input or output body.  Exact \(\Rec\) reflection
makes the shapes disjoint and gives their recognition an internal finite criterion.

For the parent collection \(P\) of the frontier admitting a derived \(o\), define
\[
  U_{G,P}(x)=\operatorname{Body}_G(x)
  \qquad\text{for every selected occurrence }x\text{ in }V.
\]
\(V\) contains exactly every assignment value, matched-reference owner, and other
stored occurrence body required by the fixed \(\Apply\) equations, with no missing
or unused member.  Each selected \(x\) lies in \(\Delta_G(P)\).  The authoritative
body equations are
\[
\begin{aligned}
 \operatorname{BaseCause}_G(o;B)
   &\Longrightarrow \operatorname{Body}_G(o)=B,\\
 \operatorname{DerivedCause}_G(o;o_{\mathrm{src}},V,\sigma,\mu,\alpha)
   &\Longrightarrow
   \operatorname{Body}_G(o)=
     \Apply(\operatorname{Body}_G(o_{\mathrm{src}}),
            U_{G,P},\sigma,\mu,\alpha).
\end{aligned}
\]
A hypothetical description application may still supply hypothetical bodies
explicitly, because those bodies are not occurrence materializations in \(G\).

\begin{theorem}[Body existence and uniqueness]\label{thm:body}
Every admitted occurrence in \(G\) has one existing, unique
\(\operatorname{Body}_G\), and its evaluation terminates.
\end{theorem}

\begin{proof}
The finite acyclic frontier graph makes strict parent ancestry well founded.  A
base body is immediate.  Every source and selected stored input of a derived cause
is in a strict parent development, so well-founded recursion supplies its unique
body.  The finite \(\Apply\) equation then supplies one unique result.  Induction
over ancestry proves the claim.
\end{proof}

The source, selected occurrence identities and exact selected positions are
causal facts.  A separately materialized result body or copied input body is not.

\subsection{Structural ancestry and immediate provenance}

Write
\[
  x<_{\mathrm{anc},G}o
\]
when \(o\) is admitted by a frontier with parents \(P\) and
\(x\in\Delta_G(P)\).  The analogous strict frontier relation is transitive parent
ancestry.  Both are well founded because the model is finite and the parent graph
is acyclic.  No numeric rank is semantic or stored.

For \(a,b\in\Delta_G(f)\), define
\[
  \operatorname{DirectDep}_{G,f}(a,b)
\]
from one actual immediate cause or endpoint position of \(a\).  A derived
occurrence depends directly on its source, every occurrence selected by the
complete/no-extra \(V\), and every exact external occurrence endpoint used by
\(\Apply\).  A base occurrence depends directly on each external occurrence
endpoint in its body.  Inverse notation may be used locally, but there is no second
dependency definition.  Transitive provenance follows the finite ancestry and
ends at base causes.

\subsection{Kernel consequences and example}

The kernel adds no implicit identity, reversal, composition, transitivity,
projection, copying, deletion, weakening, state preservation, consumption,
resource semantics, authority, or negative fact.  A place invokes no predicate
callback and asks for no declared class.  The caller supplies fillers, references,
and finite evidence; omission proves nothing.

Consider two equal-looking admitted occurrences.  If their causes are distinct base
patterns, they remain two distinct underived inputs.  If one has a derived cause,
its body is fixed by its source, complete \(V\), and \(\Apply\) maps.  Structural
equality may later compare their bodies, but it neither merges their identities nor
changes either cause.

\begin{theorem}[Renaming and stored application]\label{thm:kernel-renaming}
A one-to-one renaming of bound owner/local positions that fixes graph-backed
endpoints preserves every kernel equation.  Every derived cause projects to one
\(\Apply\) instance.  Conversely, a finite hypothetical \(\Apply\) chain may be
represented by derived causes in a Factor model once its entries and ancestry
satisfy model formation.
\end{theorem}

\section{Semantic relations}\label{sec:semantics}

\subsection{Formation, ownership, and application boundaries}

A semantic query exists only after its arguments form.  Write
\[
  \Formed_G(J)
\]
for this metalevel formation predicate.  Every proposal-defined relation below
displays a finite argument interface and a truth clause.  An extension judgment
must expose both an argument interface and one structural semantic presentation
from the next subsection.  An interface determines which tuples form; it does not
determine their truth.

For every finite supplied argument structure \(A\), let
\[
  \operatorname{Pos}(A)
\]
be exactly the positions transitively owned by \(A\).  A reference position owned
by \(A\) may end at an exact position in \(G\); that endpoint is outside
\(\operatorname{Pos}(A)\), is fixed under renaming, and is not mapped with \(A\).
A binder-owned hypothetical root and all positions transitively owned by it are
inside \(\operatorname{Pos}(A)\) and are renamed together.  This one ownership
boundary applies to every semantic argument and, later, to derivation structure.

Relative to its fixed argument interface and semantic presentation,
\(\Formed_G(J)\) holds exactly when:

\begin{enumerate}[leftmargin=2.2em]
  \item every supplied argument has the complete finite structure fixed for its
        exposed place;
  \item every graph-facing key value and external endpoint resolves uniquely in
        \(G\);
  \item every graph-backed argument occupies its required actual place and passes
        the ownership, endpoint, identity, sharing, and enclosing equations;
  \item every bound hypothetical root passes the same local equations without a
        requirement that it occur in \(G\); and
  \item all cross-argument identity, aliasing, sharing, incidence, multiplicity,
        order, and enclosing equations hold.
\end{enumerate}

Alongside formation, define the complete semantic boundary
\[
  \operatorname{AppBoundary}_G(J).
\]
It is exactly the finite structure used by the formation clause to recover the
fixed semantic definition target, every supplied ordinary argument, every
argument-place binding, and every equation needed to form the application.  It
contains no truth-evaluation footprint, proof structure, registry handle, or
implementation discriminator.  Two \(\ExactMatch\) witnesses over the same
primitive tables must project the same target, arguments, formation equations,
and formed judgment under permitted owned-position renaming; otherwise the raw
tuple is unformed.  No application decoder is an additional argument.

Graph-backed status and hypothetical status follow from
\(\operatorname{Pos}(A)\) and endpoint incidence.  No value stores a relation
name, argument type, query kind, storage-status bit, or validity bit.  Keep the
following cases distinct:

\begin{enumerate}[leftmargin=2.2em]
  \item an unformed raw tuple supplies no semantic query, positive result, negative
        result, or semantic diagnostic;
  \item a formed judgment has its stated truth value in \(G\); a signed negative
        judgment exists only when this proposal gives it its own finite formed
        interface and truth clause; and
  \item unavailable evidence leaves formation and truth unchanged and establishes
        no negation.
\end{enumerate}

\subsection{Functional structural meaning for extension roots}

An extension semantic presentation is one exact rooted position \(E\) and the
finite structure it transitively owns.  Its complete root boundary contains
exactly its argument-interface operands, truth-bearing payload operands, and their
ownership and incidence equations.  Its three candidate predicates are the
following explicit formulas:
\[
\begin{aligned}
 \operatorname{ViewShape}_G(E;A,C)
   &\Longleftrightarrow \Rec_G(E;A,C)
      \land\operatorname{ClauseGraph}_G(C,A),\\
 \operatorname{ExtShape}_G(E;A,T,e)
   &\Longleftrightarrow \Rec_G(E;A,T,e)\\[-0.3em]
   &\qquad{}\land\operatorname{TupleFamily}_G(T,A)
      \land\operatorname{ComparisonBoundary}_G(e,A,T),\\
 \operatorname{ClosureShape}_G(E;A,S,t)
   &\Longleftrightarrow \Rec_G(E;A,S,t)
      \land\operatorname{SeedFamily}_G(S,A)\\[-0.3em]
   &\qquad{}\land\exists s\,[\RefPos_G(t;s)
                         \land\operatorname{FiniteStep}_G(s,A)].
\end{aligned}
\]
For the unique ordered interface projection
\(\operatorname{Interface}_G(A;a_1,\ldots,a_n)\), define
\[
\begin{split}
 \operatorname{TupleFamily}_G(T,A)\Longleftrightarrow
 \exists X\,[\Fam_G(T;X)\land
 \forall q\in X\ \exists y_1,\ldots,y_n\,[{}&
 \operatorname{OrderedRefs}_G(q;y_1,\ldots,y_n)\\[-0.3em]
 &{}\land\bigwedge_i
 \operatorname{InterfaceBoundary}_G(a_i,y_i)]],
\end{split}
\]
where \(\operatorname{InterfaceBoundary}\) is exactly the ownership/incidence
boundary equation already used to form that argument place.  The seed-family
formula is the same equation at the unary closure interface, with \(T\) replaced
by \(S\).  The predicate \(\operatorname{FiniteStep}\)
is the already formed finite-premise application-boundary equation for its exact
target.  The predicate \(\operatorname{ComparisonBoundary}_G(e,A,T)\) holds exactly when
\(\Rec_G(e;I,X)\), \(I\) is the \(\Fam\) of all and only binder-owned positions
in \(A,T\), and \(X\) is the \(\Fam\) of exact reference positions for every
boundary-crossing endpoint; the incidence equations of both families are
reflected.  Thus \(e\) is determined by \(A,T\), has owned structure, and cannot
be the reference field \(t\) of a closure shape.

\begin{lemma}[Presentation-shape separation]\label{lem:presentation-separation}
No rooted position \(E\) satisfies two of
\(\operatorname{ViewShape}_G\), \(\operatorname{ExtShape}_G\), and
\(\operatorname{ClosureShape}_G\).  This remains true at a unary interface where
the tuple-family and seed-family equations coincide.
\end{lemma}

\begin{proof}
A view root has exactly two immediate children, whereas an extensional or closure
root has exactly three, so exact \(\Rec\) reflection separates the view case.
If the latter two predicates both held, their exact ordered root spines would
force \(T=S\) and \(e=t\).  The comparison-boundary equation requires
\(\Rec_G(e;I,X)\), so \(e\) owns exactly two positions.  The closure equation
requires \(\RefPos_G(t;s)\), which forbids \(t\) from owning any position.  This
contradiction is independent of the equations for \(T\) and \(S\).
\end{proof}

Let \(\operatorname{ClOwn}_G(C)\) be the least set containing \(C\) and closed
under \(\operatorname{Own}_G\).  The predicate
\(\operatorname{ClauseGraph}_G(C,A)\) holds exactly when there is a unique finite
pair \((N_C,r_C)\) such that:

\begin{enumerate}[leftmargin=2.2em]
  \item \(N_C\subseteq\operatorname{ClOwn}_G(C)\) and \(r_C\in N_C\);
  \item every \(n\in N_C\) satisfies exactly one reflected row below, relative
        to the fixed interface \(A\) and node set \(N_C\), and all complete
        matches of that row agree on its displayed field projections;
  \item every branch member ends in \(N_C\), and these endpoints induce the
        child-edge relation;
  \item \(r_C\) is its unique node with no incoming child edge, every node is
        reachable from \(r_C\), and the child-edge relation is acyclic; and
  \item every position of \(\operatorname{ClOwn}_G(C)\) is \(C\) or belongs to
        one of those complete row matches, and those matches preserve and reflect
        every incident primitive tuple, including every boundary-crossing tuple.
\end{enumerate}

Uniqueness here is an ordinary finite relational check.  No proposed partition,
unreachable owned structure, or alternative field projection is ignored.  To
make each row's topology literal, define
\[
\begin{aligned}
 \operatorname{Wrap}^0_G(q;x)&\Longleftrightarrow\RefPos_G(q;x),\\
 \operatorname{Wrap}^{i+1}_G(q;x)&\Longleftrightarrow
   \exists q'\,[\Rec_G(q;q')\land\operatorname{Wrap}^{i}_G(q';x)],\\
 \operatorname{Branch}^0_G(F;Y,N)&\Longleftrightarrow
   \Fam_G(F;Y)\land Y\ne\varnothing\\[-0.3em]
 &\qquad{}\land\forall u\in Y\ \exists c\,
       [\RefPos_G(u;c)\land c\in N]\\[-0.3em]
 &\qquad{}\land\forall u,v\in Y\ \forall c\,
   [\RefPos_G(u;c)\land\RefPos_G(v;c)\Rightarrow u=v],\\
 \operatorname{Branch}^{i+1}_G(F;Y,N)&\Longleftrightarrow
   \exists F'\,[\Rec_G(F;F')\land
                    \operatorname{Branch}^{i}_G(F';Y,N)].
\end{aligned}
\]
The endpoint in each member of a branch family is its child node.  Because
\(\operatorname{AppBoundary}_G(J)\) functionally recovers its one exact definition
target, write that position as \(\operatorname{Target}_G(J)\), and define the
complete formed-application operand
\[
 \operatorname{AppLeaf}_G(L;t,J)\Longleftrightarrow
 \Rec_G(L;t,J)\land\operatorname{AppBoundary}_G(J)
 \land\RefPos_G(t;\operatorname{Target}_G(J)).
\]
The ordered pair inside \(L\) is part of application formation: it exposes both
the target incidence and the entire application boundary, rather than acting as a
constructor tag.

\begin{center}
\begin{tabularx}{0.96\textwidth}{@{}>{\raggedright\arraybackslash}p{0.22\textwidth}
                                      >{\raggedright\arraybackslash}p{0.35\textwidth}
                                      X@{}}
\toprule
Clause & Exact structure & Truth contribution\\
\midrule
equality & \(\Rec(n;a,b)\), \(\operatorname{Wrap}^0(a;x)\),
\(\operatorname{Wrap}^0(b;y)\) & \(x=y\)\\
ownership & \(\Rec(n;a,b)\), \(\operatorname{Wrap}^1(a;x)\),
\(\operatorname{Wrap}^0(b;y)\) & \(\operatorname{Own}(x,y)\)\\
endpoint & \(\Rec(n;a,b)\), \(\operatorname{Wrap}^0(a;x)\),
\(\operatorname{Wrap}^1(b;y)\) & \(\operatorname{End}(x,y)\)\\
successor & \(\Rec(n;a,b)\), \(\operatorname{Wrap}^1(a;x)\),
\(\operatorname{Wrap}^1(b;y)\) & \(\operatorname{Next}(x,y)\)\\
binding & \(\Rec(n;a,b)\), \(\operatorname{Wrap}^2(a;x)\),
\(\operatorname{Wrap}^0(b;y)\) & \(\operatorname{Bind}(x,y)\)\\
formed application & \(\Rec(n;L)\),
\(\operatorname{AppLeaf}(L;t,J)\) & truth of the one formed judgment \(J\)\\
conjunction & \(\Rec(n;F)\), \(\operatorname{Branch}^0(F;Y,N_C)\)
& conjunction of the child clauses named by \(Y\)\\
projection & \(\Rec(n;q,F)\), \(\operatorname{Wrap}^0(q;c)\),
\(\operatorname{Branch}^1(F;Y,N_C)\), and
\(\exists u\in Y\ \RefPos(u;c)\) & truth of the referenced child \(c\)\\
\bottomrule
\end{tabularx}
\end{center}

\begin{lemma}[Clause-constructor separation]\label{lem:clause-separation}
For a fixed complete clause-graph candidate, no node satisfies two rows of the
eight-row table.  In particular, a projection over a singleton child family does
not also satisfy the endpoint row.
\end{lemma}

\begin{proof}
Root arity separates the two one-field rows from all six two-field rows.  If a
one-field node satisfied both remaining rows, its sole operand would both be an
\(\operatorname{AppLeaf}\), whose two owned children have an ordered
\(\operatorname{Next}\) edge, and a \(\operatorname{Branch}^0\), which forbids
every \(\operatorname{Next}\) edge between its owned children.

For the five primitive two-field rows, the wrapper-depth pairs are respectively
\[
 (0,0),\quad(1,0),\quad(0,1),\quad(1,1),\quad(2,0).
\]
Exact singleton ownership and the terminal \(\RefPos\) equation make distinct
depths incompatible.  The remaining
projection row has first depth \(0\), while its second operand satisfies
\(\operatorname{Branch}^1_G(F;Y,N_C)\): it owns one position \(F'\), and \(F'\)
owns the nonempty family \(Y\).  It therefore satisfies neither
\(\operatorname{Wrap}^0\), which owns nothing, nor
\(\operatorname{Wrap}^1\), whose sole child is a \(\RefPos\) and owns nothing.
Those are the only primitive second-operand possibilities compatible with first
depth \(0\).  If \(Y=\{u\}\), then
\(\Fam_G(F';\{u\})\) can indeed have the same local owner/next equations as
\(\Rec_G(F';u)\); nevertheless \(F'\) owns \(u\), whereas the endpoint row would
require \(F'\) itself to be a \(\RefPos\).  The extra family level therefore
separates exactly the singleton case.
\end{proof}

The wrapper, branch, and application-bundle topology is operation-defining syntax
and is included in the meaning boundary; it is not ignored metadata.  Every node,
child reference, and crossing incidence belongs to exactly one displayed clause
equation.  Consequently the clause grammar has a finite, directly checkable
criterion over
\(\operatorname{Own},\operatorname{End},\operatorname{Next},
\operatorname{Bind}\), with no callback or prose-selected interpretation.

Exactly one of the following pairwise structurally disjoint complete/no-extra
shape predicates must hold at \(E\).

\begin{enumerate}[leftmargin=2.2em]
  \item \textbf{Structural view.}  A finite clause graph exposes a conjunction or
        projection.  Each leaf is either a primitive proposal-fixed ownership,
        incidence, equality, endpoint, or binding equation, or a complete
        formed-application template for one fixed target and all arguments.
        Truth is exactly the indicated conjunction or projection of the primitive
        equations and formed leaf judgments.
  \item \textbf{Finite extensional relation.}  A complete finite tuple family
        \(T_E\) contains exactly the admitted argument tuples.  Its displayed
        comparison boundary \(e\) contains every owned position that may be
        renamed and every external endpoint that must remain fixed, with no truth
        bit.  The relation holds of \(\vec a\) exactly when \(\vec a\) equals one
        tuple by that root-preserving boundary comparison.
  \item \textbf{Fixed semantic closure.}  The payload exposes one exact finite seed
        \(S_E\) and one exact already formed finite-premise step-definition target
        \(\operatorname{Step}_{E,G}\).  Its unary relation is
        \[
          y\in\Closure_{G,\operatorname{Step}_E}(S_E).
        \]
        The step definition cannot inspect that closure.
\end{enumerate}

The patterns differ by the displayed meaning-bearing ownership and incidence,
never by a host marker, priority, kind field, or nominal constructor.  Satisfying
none or more than one makes the root malformed.  Two witnesses for one predicate
must agree on the whole boundary, target, binding, truth payload, formed judgment,
and truth clause under permitted renaming.  No residual position, edge, owner, or
closed substructure may exist only to select a pattern.

An argument-interface operand must participate in formation independently of
pattern selection.  Every truth-payload operand must participate in that pattern's
truth definition.  A judgment-template leaf instantiates the whole
\(\operatorname{AppLeaf}_G(L;t,J_\ell)\), including the whole
\(\operatorname{AppBoundary}_G(J_\ell)\) and a reference to the same fixed target;
it is not a link to a relation name.  The stored template neither proves the leaf
nor contains a truth value; the enclosing view clause invokes the already defined
truth of \(J_\ell\).  The template contains no closure expansion, quantified
family, evidence, or implementation input unless that object is itself an
ordinary formation argument.

The identity of a graph-backed presentation contains its exact owner/local root in
\(G\); isomorphic roots at different graph positions remain different relations.
A binder-owned hypothetical presentation is compared by one root-preserving
ownership/incidence isomorphism that renames all and only its internal bound
positions, transports its argument binding, preserves and reflects order,
multiplicity, identity, aliasing, sharing, and enclosing equations, and fixes every
external graph position, graph-facing key value, and graph-backed definition root.
The same equality applies to an admissibility schema root.  A root argument is
never an allocation handle, hash, byte string, or implementation address.

\subsection{Finite audit of meaning dependencies}\label{sec:audit}

Each exact extension root or admissibility schema root \(r\) has one finite
formation-and-truth audit.  Let \(\operatorname{RootCarrier}_G(r)\) contain:

\begin{itemize}[leftmargin=2em]
  \item every root in finite \(G\) satisfying one of the explicit presentation,
        schema, or finite-step shape predicates;
  \item \(r\) and every such candidate root transitively owned by \(r\); and
  \item every such root owned by a finite hypothetical structure explicitly
        supplied in the unique boundary of \(r\).
\end{itemize}

This is a finite bounding set, not an edge set or a truth quantification.
Proposal-owned definitions have two node forms.  A construction whose outgoing
meaning dependencies never depend on a supplied argument has one proposal-fixed
definition node.  An application that structurally binds a target consumed by its
dependency equation has one structural occurrence
\[
  \operatorname{DefOcc}_G(a).
\]
Here \(a\) contains the complete application-pattern incidence and the exact
dependency-bearing target-binding incidence projected from the unique
\(\ExactMatch\) witness for
\(\operatorname{AppBoundary}_G(J)\).  It retains ordinary argument-place shape but
not actual roots at data-only places.  Clause text, clause identity, and equation
names are not components of \(a\).

The target incidence preserves and reflects exact target identity, ownership,
multiplicity, aliasing, sharing, and only the order encoded by the source.  An
argument enters \(a\) exactly when the fixed structural
dependency-target projection contains its target-binding incidence.  No
counterfactual test asks whether changing that argument could change a dependency.
A position that is both an ordinary argument and dependency-bearing remains in
both projections.

At one application boundary, all clauses projecting the same \(a\) must recover
the same formed application and outgoing dependency set.  Disagreement makes the
application unformed.  Equal occurrences projected from different boundaries may
have different data-only arguments, but their outgoing dependency set must agree;
dependencies are never unioned across other occurrences.

Let \(\operatorname{OccurrenceCarrier}_G(r)\) contain only the
\(\operatorname{DefOcc}_G(a)\) values exposed by complete raw application
boundaries of \(r\) and of exact definition targets recursively exposed from those
boundaries.  This expansion is bounded by
\(\operatorname{RootCarrier}_G(r)\) and the finite proposal clauses.  Define
\[
\begin{split}
 \operatorname{AuditCarrier}_G(r)
   ={}&\operatorname{RootCarrier}_G(r)\\
      &{}\cup\{\text{proposal-fixed definition nodes}\}\\
      &{}\cup\operatorname{OccurrenceCarrier}_G(r).
\end{split}
\]

A dependency target is an extension-presentation root, an exact
admissibility-schema root, a proposal-fixed definition node, or a structural
\(\operatorname{DefOcc}_G(a)\).  The admissibility case targets the supplied
schema root, not a relation with a substitutable schema argument.  Each target
belongs to the carrier and is fixed before binder substitution.  A graph scan,
key lookup, registry, ambient relation set, or variable target cannot select one.

Add \(x\to y\) exactly when the formation or truth clause of \(x\) invokes the
formation or truth of \(y\).  Edges are projected from complete
\(\ExactMatch\) witnesses over the primitive tables
before the starting judgment is accepted as formed.  A view reaches every
judgment-template target, a semantic closure reaches its exact step target, a
property-projection occurrence reaches its exact property target, and an
admissibility schema reaches every fixed target and formation dependency of its
hypotheses and conclusion.
Every proposal-fixed node has its one proposal-stated outgoing set; every
structural occurrence has the set determined by its dependency equation and exact
target-binding incidence.

The semantic audit is the least dependency closure of \(\{r\}\) inside the exact
finite carrier.  The induced combined formation-and-truth graph must be acyclic.
Direct or mutual cycles, targets outside the carrier, variable targets, and targets
selected only after substitution make the presentation unformed.  Internal
iteration in least closure is not a dependency back edge: its already formed step
target is fixed before the closure and cannot depend back on it.  Proof edges and
retained ancestry never participate in this audit.

\begin{theorem}[Meaning determinacy]\label{thm:meaning}
The three presentation predicates are disjoint, functional, and exhaustive
on every formed presentation root.  Root equality transports every node's unique
outgoing dependency set.  The three carriers and their reachable dependency
closure are exact, finite, decidable, and invariant under permitted renaming.
Every formed extension or admissibility judgment therefore has one well-founded
formation and truth definition.
\end{theorem}

\subsection{Least closure and contexts}

For exact finite seed \(S\) and an already fixed finite-premise structural relation
\(\operatorname{Step}_G(P,y)\), define once
\[
  \Closure_{G,\operatorname{Step}}(S)
\]
as the least set containing \(S\) and closed under
\(\operatorname{Step}_G\).  The step may inspect only the finite material of one
proposed instance.  It cannot inspect this closure or a property later evaluated
over it.  Closure is semantic meaning, not an instruction to enumerate its
members.  Stage numbers may occur inside a theorem proof but are not semantic
relations.

A context at frontier \(f\) is exactly an exact finite set \(X\) satisfying
\[
  X\subseteq\Delta_G(f)
  \quad\text{and}\quad
  x\in X\land\operatorname{DirectDep}_{G,f}(x,y)
  \Longrightarrow y\in X.
\]
No constructor is required for validity.  A whole development, downward closure,
removal of an occurrence and its dependents, and finite union are ways to calculate
such a set.  Their finite descriptors are optional evidence for those calculations,
not different context meanings.  Context is an exact argument of a scoped
judgment, never a search filter or trust assumption.

\subsection{Descriptions, equality, reachability, and novelty}

A rooted description over \(G\) is a finite graph with one stored occurrence root
or one bound hypothetical root.  Bound nodes carry five-part bodies.  Endpoints may
name bound nodes or exact external positions in \(G\).  Bound names may be renamed;
external graph positions may not.  Causes, frontiers, encodings, and storage
locations are absent.

\[
  \SameStructure_G(d_1,d_2;\theta)
\]
holds when the supplied root-preserving node/position bijection \(\theta\)
preserves and reflects all five collections, ownership, direction, endpoints,
multiplicity, sharing, aliasing, and open obligations.  The root address is ignored,
but every nonroot external graph endpoint is exact.  This is isolated structural
equality, not a substitution rule for identity-sensitive application.
When the witness is omitted, use the existential abbreviation
\[
 \SameStructure_G(d_1,d_2)
 \Longleftrightarrow
 \exists\theta:\SameStructure_G(d_1,d_2;\theta).
\]

Let \(\partial_G d\) contain every exact external graph endpoint and the root
position when the root is graph-backed.  The relation
\[
  \BoundaryCorresponds_G(d_1,d_2;\theta,\beta)
\]
uses a root-preserving node/position map \(\theta\) and explicit boundary bijection
\(\beta\).  Roots are both bound or both graph-backed.  Bound endpoints follow
\(\theta\); graph endpoints follow \(\beta\).  Families of such maps used by one
application are compatible exactly when their embedded unions are partial
bijections on both occurrence references and owner/local positions.  Identity
\(\beta\) is called use equivalence only as a local abbreviation.

\begin{theorem}[Boundary-aware congruence]\label{thm:congruence}
Jointly compatible boundary correspondences for the source, inputs, and
reference-owner descriptions transport one complete description-level
\(\Apply\) instance to a corresponding instance, including assignment, match,
sharing, and fresh-position pairing.  The outputs satisfy the induced boundary
correspondence.  No such claim follows from root-erasing structural equality
alone.
\end{theorem}

For a context \(X\), let \(\operatorname{Seed}_{G}(X)\) contain the exact stored
descriptions rooted in \(X\).  Write
\[
  \operatorname{DescriptionStep}_G(P,y)
\]
for one complete description-level \(\Apply\) whose source, every input, and every
reference-owner description occurs in finite support \(P\), with every map,
multiplicity, and sharing choice exposed.
Define
\[
\begin{aligned}
  \Reach_G(X)
    &=\Closure_{G,\operatorname{DescriptionStep}}
         (\operatorname{Seed}_{G}(X)),\\
  \Derivable_{G;X}(r)
    &\Longleftrightarrow
      \exists d\in\Reach_G(X)\ \SameStructure_G(d,r),\\
  \Novel_{G;X}(r)
    &\Longleftrightarrow \neg\Derivable_{G;X}(r).
\end{aligned}
\]
A description belongs to reachability exactly when a finite valid application
trace reaches it up to renaming of bound positions.  Positive derivability is
semidecidable; novelty need not be decidable.  Failure to find a trace establishes
nothing.

For \(c\in\Delta_G(f)\), let \(X_{G,f,c}\) be the exact context obtained by removing
\(c\) and every occurrence that depends transitively on \(c\).  Current scope is
the pair \((X_{G,f,c},r)\).  Contextual scope is any supplied valid pair \((X,r)\).
Admission scope is current scope at the unique frontier admitting \(c\).  These are
abbreviations deriving one exact semantic pair, not scope kinds.  At admission,
\(X_{G,f,c}\) is the parent development; hence a derived occurrence is derivable
there and cannot be novel in admission scope.  A base cause is necessary but not
sufficient for admission-scope novelty.

\subsection{Factorization, bases, and correspondence}

Fix context \(X\).  A factorization supplies a finite composition \(L\) of
description applications from \(X\), final description \(D_L\), candidate \(D\),
and zero or one remainder.  With empty remainder it holds exactly when
\(\SameStructure_G(D,D_L)\).  With a nonempty remainder it supplies a disjoint
nonempty partition of the nodes of \(D\), one rooted expanded part containing the
root, one well-formed rooted remainder, and a complete interface table.

Every item stays with its owner.  External endpoints remain exact.  A crossing
required-link or introduced-target endpoint is redirected to one ordinary boundary
place; a crossing outgoing reference becomes one distinct introduced target to
such a place.  Endpoints to the same opposite node share a place, endpoints to
different nodes do not, and multiplicity is preserved.  The interface has exactly
one entry per cut endpoint.  The finite operation
\(\operatorname{Glue}\) is named only for this nonempty cut and reverses every
rewrite.  Factorization requires
\[
  \SameStructure_G(\operatorname{Glue}(D_e,r,I),D)
  \quad\text{and}\quad
  \SameStructure_G(D_e,D_L).
\]
Several factorizations may hold.  Boundary places are ordinary places; only the
supplied interface participates in glue.

For background context \(X\) and finite candidate set \(Q\), basis truth is the
conjunction of sufficiency and irredundancy.  Sufficiency says each exact target is
derivable from the dependency closure of \(X\cup Q\).  For \(q\in Q\), form the
dependency closure of \(X\cup(Q\setminus\{q\})\).  If it reintroduces \(q\),
irredundancy fails.  Otherwise irredundancy is exactly novelty of the rooted
remainder carried by \(q\) in that context.  No origin-based exclusion is added.

A projected correspondence fixes one finite acyclic projection schema \(\tau\)
and construction \(\pi\) before its endpoints.  Both finite traces instantiate
that same schema and their results satisfy
\(\BoundaryCorresponds_G\).  Unrelated obtainable results do not correspond.
A recovery construction is also fixed before its result and may inspect its
endpoint only through the stated correspondence.  One-way correspondence proves
one inclusion; exact transport requires a reverse construction or completeness
argument, and projection needs a lifting argument when it can merge structure.
Correspondence remains internal and grants no external interpretation.

\subsection{Generated families and recognition}\label{sec:families}

For exact context \(X\), let \(\mathcal D_{G,X}\) contain exactly the rooted
descriptions whose graph-backed root and every exact external endpoint belong to
\(X\).  Fix a finite seed \(S\subseteq\mathcal D_{G,X}\) and a finite construction
family \(\vec M\) before the seed.  A generator step
\[
  \operatorname{GeneratorStep}_{G;X,\vec M}(P,y)
\]
is one complete \(\Apply\) instance whose source is an exact indexed member of
\(\vec M\), whose premise support \(P\), result \(y\), maps, multiplicity, and
sharing are exposed, and whose every graph occurrence belongs to \(X\).  Every
premise and result lies in \(\mathcal D_{G,X}\).  Selecting a member of \(\vec M\)
does not add it to the generated family.

Use the local abbreviation
\[
 \operatorname{StructuralEqualityStep}_{G,X}(P,y)
 \Longleftrightarrow
 P=\{x\}\land x,y\in\mathcal D_{G,X}
       \land\SameStructure_G(x,y).
\]
Then
\[
  \Family_{G;X,S,\vec M}
  =
  \Closure_{G,\,
    \operatorname{GeneratorStep}_{G;X,\vec M}
      \cup\operatorname{StructuralEqualityStep}_{G,X}}(S).
\]
The explicit domain guard prevents root-erasing equality from importing a stored
root outside \(X\).  Generation is exact even when membership is undecidable.

For an already formed structural property \(K_G(d)\), define
\[
  \Necessary_{G;X,S,\vec M}(K)
  \Longleftrightarrow
  \forall d\,
  \bigl(d\in\Family_{G;X,S,\vec M}\Longrightarrow K_G(d)\bigr).
\]
The semantic meaning is universal truth over this generated family.  Finite seed,
preservation, and equality-transport derivations can prove it later but do not
define it.  Universal correspondence and recovered necessary structure are
instances with their correspondence schema or recovery construction fixed before
the quantification.

An optional recognizer is a finite family of rooted structural patterns.  Its
variables are actual open places.  Repeated uses are positions linked to the same
place.  A recursive use is an actual child-match edge whose subject is filled from
that place and which carries a positive proper-subdescription decrease.  No
variable set, recursive subset, hole role, or enable bit is stored.
\(\operatorname{Recognizes}_G(\vec N,d)\) is the least structural matching relation
induced by those places and child edges.

Recognizer adequacy has two separate semantic directions:
\[
\begin{aligned}
 \operatorname{RecognizerForward}_{G}(X,S,\vec M,\vec N)
 &\Longleftrightarrow
 \forall d\,
   (d\in\Family_{G;X,S,\vec M}
      \Rightarrow\operatorname{Recognizes}_G(\vec N,d)),\\
 \operatorname{RecognizerReverse}_{G}(X,S,\vec M,\vec N)
 &\Longleftrightarrow
 \forall d\in\mathcal D_{G,X}\,
   (\operatorname{Recognizes}_G(\vec N,d)
      \Rightarrow d\in\Family_{G;X,S,\vec M}).
\end{aligned}
\]
Their conjunction may be called exact recognition in prose.  Reverse adequacy
requires each generator premise either to have an independent finite family trace
or to be recovered as a pattern-bound smaller subdescription, preserving identity,
positions, multiplicity, and sharing.  An erased ungrounded predecessor is not
visible.  This premise-visibility condition limits this finite exactness method,
not generated-family membership.

\subsection{Structural property projection}

Preservation, transport, necessity, and invariant novelty quantify only over
properties whose meanings already exist.  Fix one exact formed application target
for a semantic relation \(B_G\), one rooted subject position \(q\) and every
position in its identity/alias class forced by ownership and sharing, every
remaining exact argument \(\vec a\), and the simultaneous replacement and
enclosing-formation equations.  Define
\[
  K_{G;B,q,\vec a}(x)
  \Longleftrightarrow
  B_G(a_1,\ldots,x,\ldots,a_n)
\]
after simultaneous replacement of exactly that alias class.

The subject may be binder-owned hypothetical structure.  Equal-looking positions
outside the forced alias class are not replaced.  Projection identity is the exact
formed tuple \((G;B,q,\vec a)\) with its aliasing and enclosing structure, modulo
permitted renaming of bound internal owners.  A relation name, executable predicate,
callback, or later evidence layout exposes no property.

The complete projection application produces one
\(\operatorname{DefOcc}_G(a)\) whose target-binding incidence contains exact target
\(B\).  Its semantic audit reaches \(B\) and the fixed projection dependencies.
Projections over different targets remain different occurrences and their
dependencies are never merged.

\subsection{Optional raw authority and bridges}\label{sec:authority}

Authority is an optional semantic protocol over exact structure in \(G\).  Its
ordinary argument \(A\) selects one initialization structure; it is not an
authority kind or namespace.  Every public authority judgment supplies its exact
frame.  A complete raw frame is
\[
  \xi=(v,f,V,Q,T;\delta_V,\delta_S),
\]
with exact version and frontier positions, proposed version view \(V\),
continuation-selection view \(Q\), selected set \(T\), and their two complete
descriptors.  Failure of one of these structural places makes a raw tuple
unformed.  If all places form but a containment or authority condition fails, the
positive predicate is false.  A negative authority judgment exists only when its
separately defined finite complement is supplied.

One exact direct view edge reads one actual link and its endpoints, without reading
a path or later predicate.  The version view is its least closure from the exact
version.  One exact continuation edge and one exact terminal-selection edge are
defined in the same direct way from \(A\)'s fixed structural selection schema.
The selection view is continuation closure; selected targets are exactly the
terminal endpoints arising from it.  Descriptors prove exact generation and
closure and include every direct endpoint.  Selection is set-valued; a required
uniqueness equation rejects a conflict instead of choosing by priority.

Use the proposal-fixed direct-closure projections
\[
\begin{aligned}
 \operatorname{ViewStep}_{A,G}(P,y)
 &\Longleftrightarrow
   \exists x\in P:\operatorname{ViewEdge}_{A,G}(x,y),\\
 \operatorname{SelectStep}_{A,G}(P,y)
 &\Longleftrightarrow
   \exists x\in P:\operatorname{SelectEdge}_{A,G}(x,y).
\end{aligned}
\]
First define the version-independent frame core
\[
\begin{split}
 \operatorname{FrameBoundary}_{A,G}(\xi)
 \Longleftrightarrow{}&
 \delta_V\text{ validates }
       V=\Closure_{G,\operatorname{ViewStep}_{A}}(\{v\})\\
 &{}\land V\subseteq\Delta_G(f)\\
 &{}\land \delta_S\text{ validates }
       \Bigl[
       Q=\Closure_{G,\operatorname{SelectStep}_{A}}(\{v\})\\
 &\hspace{11em}\land
       T=\{x:\exists a\in Q\
          \operatorname{SelectTarget}_{A,G}(a,x)\}
       \Bigr].
\end{split}
\]
Both descriptor validations are complete/no-extra and include every direct
endpoint.  Every raw core below is evaluated only after the complete frame
structure forms, includes this frame-boundary condition, and invokes neither
\(\operatorname{Version}_{A,G}\) nor a public authority wrapper.

In particular,
\(\operatorname{TransitionCore}_{A,G}(\xi,e,\xi')\) holds exactly when the
complete/no-extra transition structure at \(e\) names the exact predecessor
\(v_\xi\), an exact selected grounded method and authorized-use link, the exact
successor \(v_{\xi'}\), both frame boundaries, every explicitly changed or reused
selection, and every raw grounding or authority-core premise required by those
changes.  The successor boundary owns this generating link and admits no second
generating transition.  The core contains no proof, realization, replay,
retained-ancestry, or implementation premise.  Project it to
\[
\begin{aligned}
 \operatorname{TransitionStep}_{A,G}(P,y)
 &\Longleftrightarrow
   \exists x,e,\xi,\xi':\
   x\in P\land v_\xi=x\land v_{\xi'}=y\\[-0.2em]
 &\hspace{7em}\land
   \operatorname{TransitionCore}_{A,G}(\xi,e,\xi').
\end{aligned}
\]
The support condition \(x\in P\) is the first conjunct of the existential.
If \(v_A\) is the exact initialized version exposed by \(A\), define
\[
 \operatorname{Version}_{A,G}(v)
 \Longleftrightarrow
 v\in\Closure_{G,\operatorname{TransitionStep}_{A}}(\{v_A\}).
\]
Every noninitial version therefore exposes the actual transition that generates
it.  Neither order nor ancestry reconstructs that edge.
Now define
\[
 \operatorname{PublicFrame}_{A,G}(\xi)
 \Longleftrightarrow
 \operatorname{FrameBoundary}_{A,G}(\xi)
 \land\operatorname{Version}_{A,G}(v).
\]
Every public positive authority relation below is exactly this frame restriction
conjoined with its displayed raw core relation; no ambient frame collection is
consulted.

Grounding is one least closure from the fixed local grounding leaves and actual
support edges in the supplied frame.  A leaf is an exact initialization-audit root
or current audit-backed external binding with no grounding edge.  A nonleaf enters
when one complete local structural pattern passes and every exact support-edge
target is already in the closure.  The direct edge, not ancestry, identifies each
support.  Terminal/nonterminal names and closure stages are proof conveniences,
not semantic relations or stored fields.
The fixed leaf, local-pattern, and support-edge predicates call no transition core,
version closure, public wrapper, replay relation, or later authority result.  Thus
the exact transition step may depend on grounding without grounding depending back
on versionhood.

For an exact claim root, derive the functional projection
\[
  \operatorname{ClaimCore}_{A,G}(\xi,\mathit{claim})=(J,z).
\]
It reads only the fixed claim and target positions: \(J\) is the formed semantic
judgment asserted by the core and \(z\) is one actual owner/local root position in
\(G\).  Evidence sites, proof fragments, and retained boundaries are excluded.
Adding or moving structure outside those positions cannot change the projection.

Raw adoption is the exact structural selection
\[
  \operatorname{Adopts}_{A,G}(\xi,\mathit{claim},J,z).
\]
It holds exactly when an actual selected-claim link in the frame targets that exact
claim and its functional core is \((J,z)\).  It is extensional in the frame, exact
claim identity, and core pair.  Another claim with the same \((J,z)\) is not thereby
adopted.  Adoption says nothing about proof validity, evidence source, retained
realization, ancestry, or usability.

The remaining raw relations use the same discipline:

\begin{itemize}[leftmargin=2em]
  \item a claim, support, challenge, response, quarantine, replacement, method, or
        external binding holds only from its complete exact local pattern, selected
        endpoints, and the one grounding closure;
  \item defeat is a selected grounded challenge with no selected grounded response
        in the exact finite selected set; quarantine and its discharge use their
        own complete finite complements;
  \item a replacement points through an exact replacement link to a selected,
        grounded, supported successor claim;
  \item an active claim is selected and supported and is neither defeated,
        quarantined, nor replaced under those exact finite relations;
  \item an authorized method use names the exact selected method and exact
        authorized-use link; ordering or ancestry cannot reconstruct that link;
  \item a transition names its exact predecessor, transition step, and successor;
        a version is the least closure from initialization under those exact
        transitions;
  \item a release names its exact released structure, publication state, manifest
        fillings, and authority bindings; it stores no liveness value; and
  \item an interpretation is governed only when an exact active claim adopts its
        internal/external binding.  Structural correspondence alone never does so.
\end{itemize}

All of these are semantic relations over exact frames and positions.  They contain
no retained-replay premise, proof source, diagnostic relation, method installation,
candidate seed, implementation-reader binding, affected-result relation, or cache
rule.

A bridge has ordered source and target positions, exact source and target frames,
exact state descriptors, explicit node/position/boundary maps, and its exact
structural correspondence, coverage, recovery, and lifting premises.  The ordered
arguments determine direction; no direction field is stored.  Forward and reverse
are separate relations over swapped arguments.  One-way structure proves only its
direction.  A bridge gains authority only when an exact authority claim adopts it.
Ancestry may later establish retention but never reconstructs a bridge endpoint,
transition, release, or direction.

\begin{theorem}[Raw-authority noninterference]\label{thm:raw-authority}
If two models agree on the exact frame, selected claim, its formed core
\((J,z)\), and every raw-authority support position, then changes confined to
evidence realization, retained realization-selector links, or separately supplied
\(\operatorname{DirectUseAt}_G\) consumer structure cannot change
\(\operatorname{ClaimCore}\) or \(\operatorname{Adopts}\).
\end{theorem}

\subsection{Finite domains and universal admissibility}\label{sec:admissibility}

Define the one proposal-fixed finite guard relation
\[
  \FiniteMember_G(D,x).
\]
Its raw domain criterion is the explicit formula
\[
 \operatorname{DomainShape}_G(D;M,X)
 \Longleftrightarrow
 \Rec_G(D;M)\land\Fam_G(M;X)\land
 \forall q\in X\ \exists!x\,\RefPos_G(q;x).
\]
Define
\(\operatorname{Members}_G(D)=
  \{(q,x):\exists M,X\,[\operatorname{DomainShape}_G(D;M,X)
       \land q\in X\land\RefPos_G(q;x)]\}\), retaining the distinct member
positions \(q\) for multiplicity and any order explicitly represented inside a
member record.  There is no domain decoder.  For a permitted root \(x\), let
\(B_x\) contain
the root and every transitively owned internal position.  Its comparison boundary
also records every ownership or reference incidence touching the interior and
every exposed external endpoint and equation.

\[
  \operatorname{MemberIso}_G(D;x,y)
\]
holds by one root-preserving bijection \(B_x\to B_y\) that sends all and only
internally owned positions to internally owned positions, preserves and reflects
ownership, ordered endpoint incidence, identity, multiplicity, aliasing, sharing,
and every exposed boundary equation, and fixes every graph position outside both
interiors, every graph-facing key value, and every restricted binding result.  Its
inverse witnesses the same comparison from \(y\) to \(x\); the relation is
symmetric, not candidate-to-representative oriented.

\(D\) is formed only when its \(\operatorname{DomainShape}\) projection is
functional and no two distinct listed positions \((q,x),(q',y)\), with
\(q\ne q'\), satisfy the member-isomorphism relation on \(x,y\).  The empty
family is permitted.
Duplicates make \(D\) malformed rather than disappearing.

The judgment \(\Formed_G(\FiniteMember_G(D,x))\) holds exactly when \(D\) has one
duplicate-free finite payload and \(x\) has the required rooted boundary.  Define
\[
\begin{aligned}
 \FiniteMember_G(D,x)
 \quad\Longleftrightarrow\quad&
 \Formed_G(\FiniteMember_G(D,x))\\[-0.2em]
 &{}\land
 \exists!\,(q,y)\in\operatorname{Members}_G(D):\\[-0.2em]
 &\operatorname{MemberIso}_G(D;x,y).
\end{aligned}
\]
\(\FiniteMember\) is one proposal-fixed audit node reading raw \(D,x\); it does
not create a family of named domain predicates.

Now fix one exact finite structural schema root \(S\).  Its structural criterion is
\[
\begin{split}
 \operatorname{SchemaShape}_G(S;B,E,H,C,T,D)
 &\Longleftrightarrow \Rec_G(S;B,E,H,C,T,D)\\
 &{}\land\operatorname{BinderFamily}_G(B)
 \land\operatorname{EquationFamily}_G(E,B)\\
 &{}\land\operatorname{ApplicationFamily}_G(H,T)
 \land\operatorname{ConclusionApplication}_G(C,T)\\
 &{}\land\operatorname{FixedTargetFamily}_G(T)
 \land\operatorname{GuardFamily}_G(D,B).
\end{split}
\]
The payload names in this display are abbreviations for the following exact
record equations.  In particular, \(\operatorname{BinderFamily}_G(B)\) means
\(\Fam_G(B;X_B)\), with each \(b\in X_B\) satisfying
\(\Rec_G(b;r_b)\) for one complete finite hypothetical root \(r_b\).
\(\operatorname{EquationFamily}_G(E,B)\) means \(\Fam_G(E;X_E)\), with every
\(q\in X_E\) satisfying
\[
 \exists l,r,x,y\,[\Rec_G(q;l,r)\land\RefPos_G(l;x)
   \land\RefPos_G(r;y)\land x,y\in
   \operatorname{Pos}(\{r_b:b\in X_B\})].
\]
\(\operatorname{FixedTargetFamily}_G(T)\) means \(\Fam_G(T;X_T)\) and exactly
one target endpoint for each member.  The predicate
\(\operatorname{ApplicationFamily}_G(H,T)\)
means \(\Fam_G(H;X_H)\), with every \(a\in X_H\) satisfying
\[
 \exists t,J,z\,[\Rec_G(a;t,J)\land\RefPos_G(t;z)
   \land z\in\operatorname{Endpoints}_G(X_T)
   \land\operatorname{ApplicationRecord}_G(J;z)],
\]
where \(\operatorname{ApplicationRecord}_G(J;z)\) means that \(J\) is exactly
the whole \(\operatorname{AppBoundary}\) record formed at target \(z\), not a
relation name.  The predicate \(\operatorname{ConclusionApplication}_G(C,T)\)
is the same two-field equation for exactly one application rather than a family.
Finally, \(\operatorname{GuardFamily}_G(D,B)\) means \(\Fam_G(D;X_D)\), with
every \(g\in X_D\) satisfying
\[
\begin{aligned}
 \exists b,d,r_b,D_0\,[{}&\Rec_G(g;b,d)
 \land\RefPos_G(b;r_b)\land\RefPos_G(d;D_0)\\[-0.3em]
 &{}\land\exists q\in X_B\,\Rec_G(q;r_b)\\[-0.3em]
 &{}\land\exists M,X\,\operatorname{DomainShape}_G(D_0;M,X)].
\end{aligned}
\]
Here \(\operatorname{Endpoints}_G(X)\) is the set obtained by the unique
\(\RefPos\) projections of members of \(X\).  No payload record and no target or
guard may be unused.  These are finite formulas over the primitive tables, not an
invocation of a schema decoder.  Two witnesses must project the same binder roots,
equations, hypothesis and conclusion boundaries, fixed targets, and guards under
permitted owned-position renaming.  Define
\(\Formed_G(\UniversalAdmissible_G(S))\) to hold exactly when this unique
\(\operatorname{SchemaShape}\) witness
exists, every exposed structural place and fixed target has its required formation
boundary, every guarded domain has its duplicate-free formed payload, and the
finite semantic audit rooted at this exact schema passes.  In particular, formation
implies that there is no nonempty semantic-dependency path from the schema back to
itself.  It does not require enumerating or accepting any substitution.

A root substitution \(\theta\) maps every binder root to one finite hypothetical
image root, with no missing or extra root.  Its induced position map
\(\widehat\theta\) preserves and reflects ownership, incidence, identity,
aliasing, sharing, and every substitution equation.  Additional internal image
positions are allowed only where the exposed argument interface permits them.

\(\operatorname{FormedInstantiations}_G(S,\theta)\) is a total relation on finite
substitution candidates.  It holds exactly when substitution makes every
hypothesis and the one conclusion boundary recover one exact judgment under its
already fixed target and each recovered judgment is formed.  Missing, extra,
multiply recognized, target-changing, or unformed applications make it false.  Only
then are \(\operatorname{Hyp}_G(S,\theta)\) and
\(\operatorname{Con}_G(S,\theta)\) defined.

From the unique \(\operatorname{SchemaShape}\) witness derive the least exact projection
\[
  \operatorname{SchemaBoundary}_G(S)
    =(B_S,P_S,I_S,R_S),
  \qquad R_S=\operatorname{Bind}_G|_{I_S},
\]
by these rules:

\begin{enumerate}[leftmargin=2.2em]
  \item \(B_S\) contains exactly the binder roots, substitution equations,
        hypothesis and conclusion application boundaries, fixed
        definition/context/domain roots, guards, schema-owned domain payloads,
        formation-target projections, and all recovered incidence equations.
  \item Seed \(P_S\) with every graph-backed root, external endpoint, and
        graph-backed key place in those operands.
  \item When a formation or exact/no-extra equation is rooted in \(P_S\), add its
        owner, complete incident boundary, and all incidence endpoints.
  \item When a graph-backed definition, context, or domain root enters \(P_S\),
        add the complete operands and external endpoints from its formation
        boundary and reachable raw semantic audit.  For a finite guard, add the
        complete member payload.
  \item When a key place enters, add its value to \(I_S\), include its positive or
        absent result in \(R_S\), and add every positive endpoint to \(P_S\).
  \item Repeat the preceding closure rules to their least fixed point.
\end{enumerate}

For a complete substitution image, let
\[
  \operatorname{FullImageBoundary}_G(\theta)
    =(P_\theta,I_\theta,R_\theta)
\]
traverse every position owned by every image root, including unused internal
structure.  It records all external graph endpoints, all image key values, and
\(R_\theta=\operatorname{Bind}_G|_{I_\theta}\), including absence.

\(\operatorname{BoundaryAgreement}_G(S,\theta,\widehat\theta)\) holds exactly when
one functional correspondence \(\beta_\theta\), induced by
\(\widehat\theta\) on binder-owned positions and identity on fixed positions, maps
the complete schema-visible boundary onto the complete instantiated boundary,
preserves and reflects all ownership and incidence equations, fixed targets,
guards, key places and values, and additionally
\[
  P_\theta\subseteq P_S,\qquad
  I_\theta\subseteq I_S,\qquad
  R_\theta=R_S|_{I_\theta}.
\]
Thus even unused image structure cannot introduce a new graph reference or key.

Define, in dependency order,
\[
\begin{aligned}
 &\operatorname{PermBind}_G(S,\theta)\\
 &\quad\Longleftrightarrow
 \operatorname{RootSubstitution}_S(\theta,\widehat\theta)\\
 &\qquad{}\land
 \operatorname{BoundaryAgreement}_G(S,\theta,\widehat\theta)\\
 &\qquad{}\land
 \operatorname{FormedInstantiations}_G(S,\theta)\\
 &\qquad{}\land
 \bigwedge_{(x,D)\in\operatorname{DomainGuards}(S)}
 \begin{aligned}[t]
   [&\Formed_G(\FiniteMember_G(D,\widehat\theta(x)))\\[-0.2em]
    &{}\land\FiniteMember_G(D,\widehat\theta(x))],
 \end{aligned}
\end{aligned}
\]
and
\[
\begin{aligned}
 &\UniversalAdmissible_G(S)\\
 &\quad\Longleftrightarrow
 \forall\theta\text{ with }\operatorname{PermBind}_G(S,\theta):\\[-0.2em]
 &\hspace{4em}\Bigl[
  \bigwedge_{J_H\in\operatorname{Hyp}_G(S,\theta)}
          J_H\text{ is true in }G
  \Longrightarrow
  \operatorname{Con}_G(S,\theta)\text{ is true in }G
 \Bigr].
\end{aligned}
\]
The quantification may be infinite; this is semantic meaning, not an enumeration
instruction.  No ambient scan discovers substitutions.  A deliberately finite
domain is exposed by an ordinary guard argument.

\begin{theorem}[Finite-domain and schema boundary]\label{thm:schema}
\(\operatorname{MemberIso}_G\) is a decidable equivalence on permitted rooted
interiors, and \(\FiniteMember_G\) is functional and renaming invariant.
\(\operatorname{SchemaBoundary}_G\),
\(\operatorname{FullImageBoundary}_G\), \(\beta_\theta\), and
\(\operatorname{FormedInstantiations}_G\) are functional, finite, decidable, and
renaming invariant.  So is formation of
\(\UniversalAdmissible_G(S)\); \(\operatorname{PermBind}_G\) is false whenever an
instantiated hypothesis or conclusion is unformed, is decidable for each supplied
finite substitution, and is renaming invariant.  If two Factor models agree on the
complete schema boundary, domain payload, formation projections, incident
equations, and restricted binding results, they have the same permitted binding
family.
\end{theorem}

Consequently, extending a model only outside this complete boundary cannot create
a permitted substitution.  The admissibility judgment for \(S\) is a property of
that exact schema, not registry membership, a temporal stratum, or a fact made true
by retaining evidence.

\subsection{Semantic-root example}

Let two exact extension roots have the same ordinary argument places but different
truth payloads: one projects an ownership equation and the other a finite
extensional tuple family.  Their argument-only shapes coincide, but their complete
\(\operatorname{AppBoundary}_G\) values contain different exact definition roots
and meaning-bearing incidence.  They are distinct semantic relations.  Adding a
marker to choose between them would be malformed; their functional root predicates
already determine their meanings.

\section{Closed generic proof forms}\label{sec:proof-forms}

\subsection{The common argument boundary}

Only for an already formed judgment, derive
\[
  \operatorname{ArgBoundary}_G(J)
   =\operatorname{ArgProjection}_G
        (\operatorname{AppBoundary}_G(J)).
\]
The projection is complete/no-extra and determined only by actual incidence:

\begin{enumerate}[leftmargin=2.2em]
  \item seed it with every target reached through an ordinary argument-place
        incidence used by the formed-application equation, retaining actual
        multiplicity and endpoint order;
  \item for each seed include exactly the part of its finite rooted boundary
        exposed in \(\operatorname{AppBoundary}_G(J)\): transitively owned
        positions, ordinary argument-boundary endpoints, and all ownership,
        ordered-endpoint, identity, multiplicity, aliasing, and sharing equations
        among them;
  \item retain a position reached through both argument and definition-target
        incidence, but exclude positions and incidence reached only through a
        definition target; and
  \item take exactly this closure, rejecting every omission or extra position.
\end{enumerate}

The projection is total, functional, and renaming invariant.  It contains no
target-only discriminator.  Two supplied judgments may therefore share one
\(\operatorname{ArgBoundary}_G\) while having different fixed semantic targets.
The supplied \(J\), not proof structure, selects the requested meaning.

For finite derivation structure \(d\), use Section~\ref{sec:semantics}'s ownership
rule unchanged: \(\operatorname{Pos}(d)\) is exactly the positions transitively
owned by \(d\).  An owned reference may expose an exact endpoint in \(G\) outside
that domain; a binder-owned hypothetical position is inside.  No proof-specific
storage status is introduced.

\subsection{Complete generic catalogue}

The following proposal-owned catalogue is closed.  Form names are exposition only.
Structural applicability, not a relation-signature whitelist or stored selector,
determines which forms match.
Any \(\Formed_G(J)\), including an extension judgment, may instantiate every form
whose structural parameter and evidence equations fit it.  The catalogue is
closed in checking algorithms, not in semantic judgment signatures.

\begin{longtable}{>{\raggedright\arraybackslash}p{0.20\textwidth}
                  >{\raggedright\arraybackslash}p{0.72\textwidth}}
\toprule
Checker form & Exact finite obligation\\
\midrule
\endhead%
Direct finite &
Exactly one enumerated DF clause in the DF tables below, under
that clause's complete raw pattern.\\
Acyclic trace &
One finite rooted evidence DAG for an already defined step relation.  Each step
has its fixed finite witness or an open premise; antecedents precede their uses in
the derived topological order; no supplying proof node is embedded in local
evidence.\\
Complete enumeration &
For an explicit formed finite \(D\), exactly one representative of each
\(\operatorname{MemberIso}_G\) class denoted by its payload and none extra; every
representative passes the displayed finite obligations and the exact semantic
projection equals the claimed set.\\
Exact composite &
Every displayed semantic conjunct projects one actual open premise for its exact
formed judgment.  No local checker declares those premises true.\\
Parametric structural rule &
One finite schema with fixed local formation, substitution, and structural-decrease
checks; every semantic hypothesis is open, and an extension schema also exposes
the exact \(\UniversalAdmissible_G(S)\) premise.\\
Complete negative &
The formed semantic relation exposes one exact finite complement domain; evidence
enumerates it completely and supplies every required positive mismatch.\\
Structural pattern &
One exact rooted-pattern match with actual open places, repeated links,
child-match edges, and edge-local positive decrease checks.\\
Exact reuse &
One open premise requiring the identical formed judgment, including all arguments,
fixed endpoints, identity, and aliasing equations.\\
\bottomrule
\end{longtable}

Exact reuse is available for every \(\Formed_G(J)\).  Its sole open premise is
that identical \(J\), so it cannot establish \(J\) without an acyclic discharge
from elsewhere or an external boundary assertion.

Each form has one finite structural applicability schema and one total terminating
checker.  Its raw pattern exposes:

\begin{itemize}[leftmargin=2em]
  \item the common unlabelled conclusion-argument places;
  \item every structural parameter position;
  \item every actual open-premise socket and its prospective application boundary;
  \item every local-witness, raw key/reference, ownership, and incidence position,
        with no missing or unused position;
  \item constructive seeds and finite owner/incidence/complete-star closure rules
        for its declared inputs;
  \item every rooted subject used by an already formed property projection; and
  \item exact positive complement or coverage material where required.
\end{itemize}

No Factor value contains a form name, relation name, callback, rule identifier,
direction field, extension slot, or checker index.  A form can consume a newly
supplied semantic relation only through the ordinary structural parameter places
above.  No form concludes that \(G\) is a Factor model.

\subsection{Declared-input profiles for the primitive clauses}

Every primitive row below names one of these exact Section~\ref{sec:local-match}
projection profiles.  In all profiles, first seed \(L_\chi\) with every actual
proof-owned conclusion-argument, operand, descriptor, map, and witness position
listed in the row.  Close under owners and the row's full incidence equations.
Whenever an equation quantifies over a complete incident boundary, include its
complete star.  An external endpoint seeds \(P_\chi\); every graph-backed operand,
definition, domain, and member root also seeds \(P_\chi\).  Close \(P_\chi\) under
the same listed owner, complete-star, endpoint, and formed-application-boundary
steps.  Every reached key value enters \(I_\chi\);
\(R_\chi=\operatorname{Bind}_G|_{I_\chi}\), including absence, and every positive
binding endpoint enters \(P_\chi\).

\begin{description}[leftmargin=3.3em,style=nextline]
  \item[\(\mathsf{PA}\) --- exact atoms.]
    The seeds are the two complete supplied atoms or incidence tuples and, for
    disequality, the positive mismatch position or endpoint.
  \item[\(\mathsf{PC}\) --- finite carriers.]
    The seeds include every position and member of every supplied sequence, set,
    relation, output carrier, partition block, and coverage descriptor.  A negative
    row includes every compared member and positive mismatch witness.
  \item[\(\mathsf{PM}\) --- finite relation maps.]
    The seeds include every pair of each relation, displayed domain and codomain,
    every intermediate relation for composition, and every exact output pair.
  \item[\(\mathsf{PN}\) --- finite measures.]
    The seeds include every numeral and every complete finite collection whose
    cardinality occurs, together with all summand and comparison positions.
  \item[\(\mathsf{PI}\) --- structural isomorphism.]
    The seeds include both complete rooted interiors, their complete comparison
    boundaries, and every pair in the supplied bijection.
  \item[\(\mathsf{PP}\) --- application.]
    The seeds include the source and input bodies, assignment, reference-owner
    positions, link match, fresh-position pairing, and the exact result body.
\end{description}

These are constructive projection definitions, not descriptions of runtime reads.
The fixed checker and its equations are not members of \(\chi\).

\subsection{DF-1: exact atoms}\label{sec:df1}

\begin{longtable}{>{\raggedright\arraybackslash}p{0.16\textwidth}
                  >{\raggedright\arraybackslash}p{0.61\textwidth}
                  >{\raggedright\arraybackslash}p{0.13\textwidth}}
\caption{Complete DF-1 catalogue}\label{tab:df1}\\
\toprule
Clause & Complete raw conclusion, operands, and witness & Inputs\\
\midrule
\endfirsthead%
\toprule
Clause & Complete raw conclusion, operands, and witness & Inputs\\
\midrule
\endhead%
DF1-U-EQ & \(u=v\) for two supplied values of \(U\); the two operand places are
complete. & \(\mathsf{PA}\)\\
DF1-U-NE & \(u\ne v\); the witness is their exact positive value mismatch. &
\(\mathsf{PA}\)\\
DF1-POS-EQ & equality of two supplied owner/local positions, including owners and
local coordinates. & \(\mathsf{PA}\)\\
DF1-POS-NE & disequality of owner/local positions, with the differing owner or
local coordinate exposed. & \(\mathsf{PA}\)\\
DF1-END-EQ & equality of two exact supplied endpoints, including graph-backed or
bound-root status. & \(\mathsf{PA}\)\\
DF1-END-NE & endpoint disequality, with the differing endpoint component exposed.
& \(\mathsf{PA}\)\\
DF1-INC-EQ & equality of two supplied ordered incidence tuples, position by
position and with equal length. & \(\mathsf{PA}\)\\
DF1-INC-NE & ordered-incidence disequality, with a length mismatch or one exact
differing indexed component. & \(\mathsf{PA}\)\\
\bottomrule
\end{longtable}

\subsection{DF-2: finite extensional carriers}\label{sec:df2}

A finite sequence is compared as its complete ordered index/value relation.  A
duplicate-free set is compared extensionally.  A finite relation is compared as
its complete pair set.  The following thirty rows are distinct clauses.

\begin{longtable}{>{\raggedright\arraybackslash}p{0.17\textwidth}
                  >{\raggedright\arraybackslash}p{0.60\textwidth}
                  >{\raggedright\arraybackslash}p{0.12\textwidth}}
\caption{Complete DF-2 catalogue}\label{tab:df2}\\
\toprule
Clause & Complete raw conclusion, operands, and witness & Inputs\\
\midrule
\endfirsthead%
\toprule
Clause & Complete raw conclusion, operands, and witness & Inputs\\
\midrule
\endhead%
DF2-MEM-SEQ & One supplied index/value pair occurs in the complete finite
sequence descriptor. & \(\mathsf{PC}\)\\
DF2-MEM-SET & One supplied value occurs in the complete duplicate-free set. &
\(\mathsf{PC}\)\\
DF2-MEM-REL & One supplied ordered pair occurs in the complete finite relation. &
\(\mathsf{PC}\)\\
DF2-NMEM-SEQ & The complete sequence is supplied and the candidate pair has a
positive mismatch against every indexed pair. & \(\mathsf{PC}\)\\
DF2-NMEM-SET & The complete set is supplied and the value has a positive mismatch
against every member. & \(\mathsf{PC}\)\\
DF2-NMEM-REL & The complete relation is supplied and the pair has a positive
mismatch against every pair. & \(\mathsf{PC}\)\\
DF2-EQ-SEQ & Two complete sequences have the same length and equal indexed values.
& \(\mathsf{PC}\)\\
DF2-EQ-SET & Two complete duplicate-free sets contain exactly the same values. &
\(\mathsf{PC}\)\\
DF2-EQ-REL & Two complete finite relations contain exactly the same ordered pairs.
& \(\mathsf{PC}\)\\
DF2-SUB-SEQ & Every index/value pair of the first complete sequence occurs in the
second at that same index. & \(\mathsf{PC}\)\\
DF2-SUB-SET & Every member of the first complete set occurs in the second. &
\(\mathsf{PC}\)\\
DF2-SUB-REL & Every pair of the first complete relation occurs in the second. &
\(\mathsf{PC}\)\\
DF2-DISJ-SEQ & The complete index/value pair sets of two sequences are disjoint. &
\(\mathsf{PC}\)\\
DF2-DISJ-SET & Two complete sets have no common value, with all pairwise
disequalities exposed. & \(\mathsf{PC}\)\\
DF2-DISJ-REL & Two complete relations have no common ordered pair. &
\(\mathsf{PC}\)\\
DF2-UNION-SEQ & The output sequence's complete index/value relation is exactly the
compatible union of the two input relations; shared indices have equal values. &
\(\mathsf{PC}\)\\
DF2-UNION-SET & The output set is exactly the union of two complete sets. &
\(\mathsf{PC}\)\\
DF2-UNION-REL & The output relation is exactly the union of two complete
relations. & \(\mathsf{PC}\)\\
DF2-INTER-SEQ & The output index/value relation is exactly the intersection of two
complete sequence relations. & \(\mathsf{PC}\)\\
DF2-INTER-SET & The output set is exactly the intersection of two complete sets.
& \(\mathsf{PC}\)\\
DF2-INTER-REL & The output relation is exactly the intersection of two complete
relations. & \(\mathsf{PC}\)\\
DF2-DIFF-SEQ & The output index/value relation is exactly the first complete
sequence relation minus the second. & \(\mathsf{PC}\)\\
DF2-DIFF-SET & The output set is exactly the first complete set minus the second.
& \(\mathsf{PC}\)\\
DF2-DIFF-REL & The output relation is exactly the first complete relation minus
the second. & \(\mathsf{PC}\)\\
DF2-PART-SEQ & The supplied complete block family is pairwise disjoint and its
union is exactly the sequence's index/value relation. & \(\mathsf{PC}\)\\
DF2-PART-SET & The supplied complete block family is pairwise disjoint and its
union is exactly the set. & \(\mathsf{PC}\)\\
DF2-PART-REL & The supplied complete block family is pairwise disjoint and its
union is exactly the relation. & \(\mathsf{PC}\)\\
DF2-COVER-SEQ & Every and only pair in the complete sequence descriptor occurs in
the supplied covering family. & \(\mathsf{PC}\)\\
DF2-COVER-SET & Every and only member of the complete set occurs in the supplied
covering family. & \(\mathsf{PC}\)\\
DF2-COVER-REL & Every and only pair of the complete relation occurs in the
supplied covering family. & \(\mathsf{PC}\)\\
\bottomrule
\end{longtable}

\subsection{DF-3: finite relation maps}\label{sec:df3}

\begin{longtable}{>{\raggedright\arraybackslash}p{0.18\textwidth}
                  >{\raggedright\arraybackslash}p{0.59\textwidth}
                  >{\raggedright\arraybackslash}p{0.12\textwidth}}
\caption{Complete DF-3 catalogue}\label{tab:df3}\\
\toprule
Clause & Complete raw conclusion, operands, and witness & Inputs\\
\midrule
\endfirsthead%
\toprule
Clause & Complete raw conclusion, operands, and witness & Inputs\\
\midrule
\endhead%
DF3-DOM & The supplied set is exactly the first-component domain of every pair in
the complete relation. & \(\mathsf{PM}\)\\
DF3-RANGE & The supplied set is exactly the second-component range of every pair.
& \(\mathsf{PM}\)\\
DF3-RESTRICT & The output is exactly the supplied relation restricted to the
displayed domain. & \(\mathsf{PM}\)\\
DF3-ID & The supplied relation contains exactly \((x,x)\) for each member of the
displayed finite domain. & \(\mathsf{PM}\)\\
DF3-INVERSE & The output contains exactly the coordinate reversals of every pair
in the input relation. & \(\mathsf{PM}\)\\
DF3-COMP & The output contains exactly each \((x,z)\) having an exposed
intermediate \(y\) in the two complete input relations. & \(\mathsf{PM}\)\\
DF3-FUNCTION & No two pairs with equal first component have unequal second
components; all pairs are supplied. & \(\mathsf{PM}\)\\
DF3-TOTAL & The relation is functional and its exact domain equals the displayed
finite domain. & \(\mathsf{PM}\)\\
DF3-INJECTIVE & Equal second components imply equal first components in the
complete relation. & \(\mathsf{PM}\)\\
DF3-SURJECTIVE & The complete range equals the displayed finite codomain. &
\(\mathsf{PM}\)\\
DF3-PBIJ & The relation is functional and injective between its exact displayed
domain and range. & \(\mathsf{PM}\)\\
DF3-BIJ & The relation is functional, displayed-domain total, injective, and
surjective onto the displayed codomain. & \(\mathsf{PM}\)\\
\bottomrule
\end{longtable}

\subsection{DF-4: finite measures}\label{sec:df4}

\begin{longtable}{>{\raggedright\arraybackslash}p{0.18\textwidth}
                  >{\raggedright\arraybackslash}p{0.59\textwidth}
                  >{\raggedright\arraybackslash}p{0.12\textwidth}}
\caption{Complete DF-4 catalogue}\label{tab:df4}\\
\toprule
Clause & Complete raw conclusion, operands, and witness & Inputs\\
\midrule
\endfirsthead%
\toprule
Clause & Complete raw conclusion, operands, and witness & Inputs\\
\midrule
\endhead%
DF4-NUMERAL & A supplied natural term is exactly the supplied numeral. &
\(\mathsf{PN}\)\\
DF4-CARD & A supplied natural term equals the cardinality of one complete finite
sequence, set, or relation. & \(\mathsf{PN}\)\\
DF4-SUM & A supplied natural term is the finite sum of the complete supplied term
sequence. & \(\mathsf{PN}\)\\
DF4-EQ & Two terms built only by DF4-NUMERAL, DF4-CARD, and DF4-SUM have equal
natural values. & \(\mathsf{PN}\)\\
DF4-LE & Two such terms satisfy \(\leq\). & \(\mathsf{PN}\)\\
DF4-LT & Two such terms satisfy \(<\), including structural-size decrease. &
\(\mathsf{PN}\)\\
\bottomrule
\end{longtable}

\subsection{DF-5 and DF-6}\label{sec:df56}

\begin{longtable}{>{\raggedright\arraybackslash}p{0.18\textwidth}
                  >{\raggedright\arraybackslash}p{0.59\textwidth}
                  >{\raggedright\arraybackslash}p{0.12\textwidth}}
\caption{Complete DF-5 and DF-6 catalogue}\label{tab:df6}\\
\toprule
Clause & Complete raw conclusion, operands, and witness & Inputs\\
\midrule
\endfirsthead%
\toprule
Clause & Complete raw conclusion, operands, and witness & Inputs\\
\midrule
\endhead%
DF5-ISO & One supplied finite root-preserving position bijection has exact domain
and range and preserves and reflects ownership, ordered endpoint incidence,
identity, multiplicity, aliasing, sharing, and fixed external endpoints.  This is
the sole proof-side isomorphism primitive. & \(\mathsf{PI}\)\\
DF6-APPLY & One complete five-argument \(\Apply\) instance, with its source,
inputs, reference owners, assignment, link match, fresh-position pairing, and
exact five-collection result. & \(\mathsf{PP}\)\\
\bottomrule
\end{longtable}

The catalogue therefore has exactly \(8+30+12+6+1+1=58\) primitive clauses.
The DF labels authorize no unlisted operation.  A local admission view, authority
predicate, projection, constructor equation, or arithmetic operator is not direct
merely because it is finite.  It must reduce through its semantic definition to
one listed row, fit another generic form, or occur as an open premise.  Adding a
primitive row requires a later proposal.

\subsection{Input projections for the remaining forms}

The remaining form schemas use the common closure rules above and the following
complete seeds.

\begin{longtable}{>{\raggedright\arraybackslash}p{0.20\textwidth}
                  >{\raggedright\arraybackslash}p{0.72\textwidth}}
\toprule
Form & Exact \(L_\chi\) seeds and additional closure\\
\midrule
\endhead%
Acyclic trace &
Every trace node, step-witness position, premise occurrence, root/result boundary,
and local order edge.  Close over each supplied step's complete argument boundary;
graph-backed step targets and endpoints enter \(P_\chi\).\\
Complete enumeration &
The whole proof-owned descriptor, every representative interior and boundary,
every class-comparison map, and the exact projected output set.  A graph-backed
domain or member root enters \(P_\chi\).\\
Exact composite &
Every conjunct application boundary and every actual open socket, with its
prospective argument tuple and fixed target.  Close over the complete exposed
argument boundaries, not the truth definitions.\\
Parametric structural rule &
The complete schema evidence, binder roots, substitution positions, every local
evidence node, every open socket, all substitution and strict-decrease maps, and
the conclusion boundary.  Graph-backed schema, target, domain, and argument roots
enter \(P_\chi\).\\
Complete negative &
The complete finite complement-domain descriptor, every candidate, and every
positive mismatch position.  Include the entire graph-backed domain payload in
\(P_\chi\) when applicable.\\
Structural pattern &
The whole rooted pattern image, each open place and repeated link, every child
edge, child subject, and positive proper-subdescription witness.\\
Exact reuse &
The conclusion boundary, one open socket, and the socket's complete identical
prospective application boundary.\\
\bottomrule
\end{longtable}

For each form, the proposal-fixed raw schema also lists all key/reference places.
Those values and results enter \(I_\chi,R_\chi\) by the common rule.  There is no
equation-list component, checker identity, or observed-read component.

\subsection{Extensional bindings and images}

For one raw generic-form pattern, a local binding \(m\) is the finite relation of
pairs \((u,v)\) from its proposal-displayed owner/local positions \(u\) to actual
positions \(v\) in a supplied derivation.  It preserves and reflects ownership,
incidence, equality, aliasing, and sharing.  Its domain is the finite pattern
structure modulo permitted renaming, not a schema or pattern-object identity.
Two bindings are equal exactly when an incidence-preserving isomorphism of their
pattern sides makes their actual position-pair relations equal while fixing exact
semantic endpoints in \(G\).

\(\operatorname{Image}_d(m)\) is only the actual finite positions selected in
\(d\), with ownership and incidence induced from \(d\).  It contains no
pattern-side position and no incidence object separate from its owned positions.
Every image contains the complete common conclusion boundary.

For fixed formed \(J\), each raw match projects:

\begin{itemize}[leftmargin=2em]
  \item the common conclusion-argument places relative to
        \(\operatorname{ArgBoundary}_G(J)\);
  \item the four declared-input components, including every local witness position
        but no checker identity; and
  \item every open socket, prospective premise tuple, and fixed intended semantic
        target.
\end{itemize}

It contains no discharge choice and no assumption.  For fixed
\((G,J,d,k,w)\), two complete raw matches with equal
\(\operatorname{Image}_d(m)\) must project the same arguments, prospective premise
tuples and targets, and coherent declared inputs: \(L_\chi\) agrees by permitted
induced-structure isomorphism, while \(P_\chi,I_\chi,R_\chi\) agree exactly.
Different fixed checkers may apply.  Several coherent bindings and clauses may
share one image and remain disjunctive.  No coherence condition compares different
supplied judgments.

\subsection{Parametric structural-rule format}

The parametric form is evidence inside this fixed catalogue, not a way to install
proof semantics.  One applicable instance exposes:

\begin{itemize}[leftmargin=2em]
  \item any exact property projection \(K_{G;B,q,\vec a}\), including its forced
        alias class and remaining arguments;
  \item every actual open premise and repeated-use link;
  \item the requested output judgment at the common conclusion boundary;
  \item every fixed helper judgment as an ordinary open premise;
  \item for an extension schema \(S\), the exact formed judgment
        \(\UniversalAdmissible_G(S)\) as another open premise;
  \item every template-local evidence node and incidence edge;
  \item every fixed finite graph-formation, substitution, local-check, and
        structural-decrease obligation; and
  \item the topological order derived from the template-local evidence graph.
\end{itemize}

The total checker verifies that the finite parametric DAG is valid structure, maps
retain functional keys, sets retain distinct members, endpoints resolve, semantic
hypotheses occupy actual open sockets, binder-derived owner domains are disjoint,
and undeclared identity, aliasing, sharing, or owner collision is absent.  It
contains no supplying proof node; supplying nodes belong only to the complete
derivation in Section~\ref{sec:derivations}.

If the schema contains a local recursive use, that use has an explicit dependency
edge and an incidence-preserving injection from the dependent use's transitively
owned positions into a strict subset of the source use's positions.  This strict
finite owned-substructure decrease is the well-founded condition.  No registration
time, retention order, or nominal earlier-rule order participates.

The proposal proves one substitution theorem for this generic form: truth of
\(\UniversalAdmissible_G(S)\) and every instantiated semantic hypothesis, together
with the finite local conditions, implies the instantiated conclusion.  The
checker validates one supplied finite substitution and never enumerates the
schema's permitted binding family or invokes the conclusion relation to decide
whether the template passes.

Distinct open places imply neither equality nor disequality.  Repeated identity or
sharing exists only through actual links to the same place.  Any additional
equality or disequality must be an exact structural antecedent or an ordinary open
premise.

\subsection{Complete decomposition tables}

The following are the only step and transport decomposition tables.  Row names are
exposition; the complete ordinary positions and incidence shown by each row select
it.  Within these two tables only, \(\mathsf{DS}\), \(\mathsf{GS}\),
\(\mathsf{ES}\), and \(\mathsf{FS}\) abbreviate the description, generator,
structural-equality, and combined family step relations.  Likewise
\(\mathsf{SS}\), \(\mathsf{BC}\), \(\mathsf{UE}\), \(\mathsf{FC}\), and
\(\mathsf{PC}\) abbreviate structural equality, boundary correspondence,
identity-boundary correspondence, full correspondence, and projected
correspondence.  These abbreviations add no relation.

\begin{longtable}{>{\raggedright\arraybackslash}p{0.20\textwidth}
                  >{\raggedright\arraybackslash}p{0.30\textwidth}
                  >{\raggedright\arraybackslash}p{0.42\textwidth}}
\caption{Complete DEC-S table}\label{tab:decs}\\
\toprule
Row & Exact step relation & Complete decomposition\\
\midrule
\endfirsthead%
\toprule
Row & Exact step relation & Complete decomposition\\
\midrule
\endhead%
DEC-S-DESCRIPTION &
\(\mathsf{DS}_G(P,y)\) &
One complete description \(\Apply\) and join; source, every input and
reference-owner description, exact support \(P\), assignment, match, fresh map,
multiplicity, sharing, and result \(y\) are exposed.\\
DEC-S-GENERATOR &
\(\mathsf{GS}_{G;X,\vec M}(P,y)\) &
One exact indexed source member of \(\vec M\), one complete description
\(\Apply\), exactly support \(P\), result \(y\), and all context/domain guards and
maps are exposed.\\
DEC-S-EQUALITY &
\(\mathsf{ES}_{G,X}(P,y)\) &
Exactly \(P=\{x\}\), \(x,y\in\mathcal D_{G,X}\), and one complete
\(\SameStructure_G(x,y)\) map.\\
DEC-S-FAMILY &
\(\mathsf{FS}_{G;X,\vec M}(P,y)\) &
The complete disjoint structural union of the preceding generator and equality
raw patterns.  Every matching alternative is included; no flag selects one.\\
\bottomrule
\end{longtable}

\begin{longtable}{>{\raggedright\arraybackslash}p{0.20\textwidth}
                  >{\raggedright\arraybackslash}p{0.30\textwidth}
                  >{\raggedright\arraybackslash}p{0.42\textwidth}}
\caption{Complete DEC-R table}\label{tab:decr}\\
\toprule
Row & Exact transport relation & Complete decomposition\\
\midrule
\endfirsthead%
\toprule
Row & Exact transport relation & Complete decomposition\\
\midrule
\endhead%
DEC-R-SAME &
\(\mathsf{SS}_G(x,x')\) &
Both complete rooted interiors and one exact node/position bijection preserving
and reflecting every structural equation.\\
DEC-R-BOUNDARY &
\(\mathsf{BC}_G(x,x';\theta,\beta)\) &
Both complete boundaries, node/position map, external boundary map, and every
joint compatibility equation used by the requested instance.\\
DEC-R-USE &
\(\mathsf{UE}_G(x,x';\theta)\) &
The complete DEC-R-BOUNDARY structure plus the exact equation that every external
boundary endpoint is fixed.\\
DEC-R-FULL &
\(\mathsf{FC}_G(x,x';\theta,\beta)\) &
The complete boundary-correspondence operands projected by the full
correspondence application, with no additional mode position.\\
DEC-R-PROJECTED &
\(\mathsf{PC}_{G;X,\pi,\tau}(x,x')\) &
The fixed projection schema, both complete instantiation traces, projected
endpoints, node/position/boundary maps, multiplicity, sharing, and result
correspondence.\\
\bottomrule
\end{longtable}

The identity of a decomposition match is its structural-isomorphism class.
Graph-backed owners, occurrence identities, external endpoints, positions,
multiplicity, aliasing, sharing, and constructor order remain exact.  Only
decomposition-local bound owners may be bijectively renamed.  No byte, hash,
opaque-owner order, traversal order, or least representative is used.

A complete-enumeration descriptor fixes exact formed \(D\) before proof and
contains one arbitrary structurally valid representative of every
\(\operatorname{MemberIso}_G\) class denoted by \(D\), and none extra.  Finite
DF2 and DF5 checks establish coverage without open premises merely for replaying
those checks.  Distinct decomposition classes may project the same semantic
\((P,y)\); exact equality is required only after projecting the complete descriptor
to the semantic pair set.

\subsection{Recognition's separate pattern format}

Recognizer evidence uses the semantic pattern format from
Section~\ref{sec:families}, not the parametric-rule template:

\begin{itemize}[leftmargin=2em]
  \item one fixed rooted shape with actual open places;
  \item repeated positions linked to the same place;
  \item actual child-match edges filled from those places; and
  \item one positive proper-subdescription check at each recursive edge.
\end{itemize}

The direct form checks one match.  Separate forward and reverse forms check their
finite coverage and reconstruction obligations.  Their conjunction has no third
form.  An extension cannot add a recognizer checker, proof primitive, callback, or
ambient admissible-rule registry.

Every preservation, transport, coverage, necessity, recognition direction, bridge
result, or schema admissibility result can later be used only through an ordinary
premise edge or boundary assertion.  Exact reuse imposes the same discharge rule.
A retained realization never extends this catalogue.  Complete negative applies
only to a semantic relation with a fixed finite complement and positive mismatch
evidence; failed search, an absent match, or generic validation failure never
qualifies.

\begin{theorem}[Catalogue closure and erasure]\label{thm:catalogue}
For each row and generic form, structural applicability, the local checker, and its
declared-input projection are total and terminating on supplied finite structure.
The catalogue contains exactly the forms and primitive rows displayed above.
Erasing names, serialization labels, host types, and dispatch metadata preserves
\(\operatorname{AppBoundary}_G\), \(\operatorname{ArgBoundary}_G\), every raw
match, and every declared input.
\end{theorem}

\section{Local matching in one finite derivation}\label{sec:local-match}

\subsection{Complete raw matching}

Let \(d\) be supplied finite Factor structure, \(k\) an actual local root, and
\(w\) the exact ownership/incidence-delimited neighborhood derived at \(k\).
\(w\) is not stored.  The layer receives an already formed requested judgment
\(J\), but raw matching performs no graph lookup, semantic formation of a
prospective premise, or resolved conclusion comparison.

For exposition, write
\[
  \operatorname{LocalMatch}_{F}(J,d,k;w,m)
\]
when one fixed catalogue clause's raw pattern matches every required
proof-owned conclusion-argument place, raw key/reference place,
prospective-premise socket, local-witness place, and ownership/incidence equation,
with no missing or extra pair.  A missing local place means that clause has no raw
match.  An exact external endpoint already present as an ordinary argument of
\(J\) may be compared as an atom, but no binding lookup or graph traversal occurs.

This notation is clause-local exposition, not a stored relation or selected-form
argument.  Matching is relational: several complete \(m\) values and several
clauses may match the same actual image.  Each \(m\) is one complete extensional
binding as defined in Section~\ref{sec:proof-forms}.

\subsection{Declared inputs before resolved checks}

For each complete raw match, derive
\[
\begin{split}
 \operatorname{DeclaredInputs}_{F,G}(J,d,k,w,m)
   &=\chi,\\
 \chi&=(L_\chi,P_\chi,I_\chi,R_\chi),\\
 R_\chi&=\operatorname{Bind}_G|_{I_\chi}.
\end{split}
\]
The applicable row's profile and the remaining-form table in
Section~\ref{sec:proof-forms} give constructive seeds and closure:

\begin{enumerate}[leftmargin=2.2em]
  \item seed \(L_\chi\) with the actual positions paired to conclusion, premise,
        local-witness, finite-domain, operand, map, and supplied-equation places;
  \item close it under the clause's owner and incidence rules, including the
        complete incident star wherever an exact/no-extra equation quantifies over
        that star;
  \item seed \(P_\chi\) with every graph-backed argument, endpoint, definition,
        context, domain, and member projected from \(J\) or a matched reference;
  \item close it under the listed finite owner, complete-star, endpoint, and
        formed-application-boundary rules;
  \item put every reached key value in \(I_\chi\), retain its positive or absent
        result in \(R_\chi\), and add every positive binding endpoint to
        \(P_\chi\); and
  \item take the least simultaneous closure.
\end{enumerate}

Fixed finite \(G,d,m\), and clause make this bundle finite and unique.  A
proof-owned descriptor, local member, reference, or witness is in \(L_\chi\); an
exact graph domain/member position or endpoint is in \(P_\chi\); a key and its
positive or absent result are in \(I_\chi,R_\chi\).  A complete enumeration covers
exactly the \(\operatorname{MemberIso}_G\) classes already denoted by its fixed
domain.  It cannot infer a domain by scanning \(G\), invent an order, or copy a
graph position into local proof evidence.

\subsection{Candidate and resolved match relations}

Define
\[
  \CandidateAt_G(J,d,k;w,m,\chi)
\]
as the finite disjunction of all fixed catalogue clauses, each requiring one
complete raw local match and its uniquely derived four-component bundle.  The
projection does not filter a completed raw match.  No selected form, clause, or
candidate value witnesses the disjunction.  Neither \(d\) nor a later selector
stores \(J,w,m,\chi\), or a clause.

Define
\[
  \MatchAt_G(J,d,k;w,m,\chi)
\]
by adjoining all resolved conditions to that candidate:

\begin{enumerate}[leftmargin=2.2em]
  \item \(\Formed_G(J)\) and exact structural equality of the actual common
        conclusion boundary \(C_k\) with \(\operatorname{ArgBoundary}_G(J)\);
  \item graph-relative formation and applicability of every fixed semantic
        parameter;
  \item exact resolution of every declared key and reference comparison;
  \item recovery of one exact \(\Formed_G(J_r)\) from every prospective-premise
        application, and equality of its socket boundary with
        \(\operatorname{ArgBoundary}_G(J_r)\); and
  \item success of the proposal-authored total finite checker over the declared
        structural inputs and positive witness positions.
\end{enumerate}

A boundary mismatch, malformed external or premise application, absent or wrong
lookup, failed applicability equation, or failed checker leaves
\(\CandidateAt_G\), its raw \(m\), and all of \(\chi\) intact.  Every structural
input that may vary between Factor models is declared before the result.  A fixed
checker cannot call an ambient oracle, registry, callback, extension evaluator,
unbounded search, or the semantic conclusion relation.

Only a passing resolved match derives
\[
  \operatorname{PremiseNeed}_G(m,r;J_r),
\]
where \(r\) is the actual socket and \(J_r\) the exact formed judgment recovered
there.  Two equal judgments at different sockets remain two needs.  Needs from
alternative matches are never unioned.  Local matching does not choose how a need
is discharged.

\begin{theorem}[Raw-match determinacy and failure preservation]%
\label{thm:local-determinacy}
Declared inputs are functional for each complete raw match; raw matching itself
need not be functional.  For fixed \(J\), equal actual images have coherent
arguments, prospective premises, targets, and four-component bundles.  Candidate
and resolved-match truth are determined only by the complete raw match, fixed
clause equations, and declared bundle.  Every raw candidate and bundle persists
through resolved failure.
\end{theorem}

\begin{theorem}[Locality of resolved matching]\label{thm:match-locality}
For every clause with a raw match at a selected image, if Factor models \(G,G'\)
agree on the complete raw formation inputs, all four declared components, complete
incident stars, exact endpoints, and positive or absent restricted binding results,
then the raw match, its declared bundle, every prospective-premise projection, and
the disjunction of passing resolved matches agree in \(G,G'\).  Schema-dependent
forms use the complete-image invariance of
\(\operatorname{PermBind}\).
\end{theorem}

\section{Complete finite derivations}\label{sec:derivations}

\subsection{Derivation structure and actual-side selection}

A finite derivation \(d\) contains only positions, ownership and incidence, one
common unlabelled conclusion boundary at each local root, local evidence, open
premise sockets with that same boundary format, proof-supplying edges, and
boundary-assertion witnesses.  These are structural views, not stored kinds.
\(d\) stores no local judgment, neighborhood, raw match, declared bundle,
selection, or assumption set.  Exact endpoints in \(G\) referenced by its positions
are outside \(\operatorname{Pos}(d)\).

Supply separately one finite extensional actual-side selection
\[
  h\subseteq
  \{\text{actual local roots of }d\}\times\operatorname{Pos}(d).
\]
For local root \(k\), write \(h_k=\{q:(k,q)\in h\}\).  It is one exact
actual-side substructure with ownership, incidence, and boundary inherited from
\(d\).  It includes each local evidence edge, socket, or link position selected by
the generic pattern, but excludes the proof-supplying edge or boundary-witness
position that later discharges an open socket.

\(h_k\) selects a local match exactly when some catalogue binding \(m_k\) has
\(\operatorname{Image}_d(m_k)=h_k\).  Pattern positions are never members of
\(h\).  For fixed required \(J_k\), static image coherence determines the
arguments, prospective premises, intended targets, and four-component bundle,
while allowing several coherent clauses with different fixed checkers.  The image
does not determine \(J_k\), discharge a premise, or project an assumption.

Equality of \(h\) fixes this relation on actual positions modulo one permitted
renaming of owners wholly internal to \(d\).  Form identifiers, schema identities,
pattern objects, serialization, and check order are absent.  The complete
derivation witness for supplied \(J\) is \((d,h)\); neither component alone is
complete.

\subsection{Complete matching of one reading}

Define
\[
  \CompleteMatch_G(J,d,h).
\]
At the unique root \(k_0\), set \(J_{k_0}=J\).  At every nonroot node,
\(J_k\) is derived from the exact prospective premise at its incoming proof edge.
All incoming edges to a shared node must derive the same formed judgment.

The relation holds exactly when:

\begin{enumerate}[leftmargin=2.2em]
  \item \(d\) has one structural root \(k_0\).  For some
        \(w_0,m_0,\chi_0\), the resolved-match relation passes at
        \((J,d,k_0)\), \(h_{k_0}\) is exactly the image of \(m_0\), and
        \(C_{k_0}\) is structurally equal to
        \(\operatorname{ArgBoundary}_G(J)\);
  \item every node reached from \(k_0\) through proof edges at selected sockets
        has exactly one restriction \(h_k\) that is the image of some \(m_k\)
        whose resolved-match relation passes for \((J_k,d,k)\), and every pair in
        \(h\) belongs to one such reached node and image;
  \item every selected premise need at socket \(r\), binding \(m_k\), and formed
        judgment \(J_r\) has exactly one of two structurally disjoint actual
        discharge shapes:
        \begin{enumerate}[label=\alph*.,leftmargin=2em]
          \item one proof edge to a node whose selected image passes
                \(\MatchAt_G\) for the identical \(J_r\); or
          \item one complete boundary-assertion witness whose exact projection is
                formed \(J_r\), asserting but not proving it;
        \end{enumerate}
        and both the socket and supplying conclusion boundaries equal the exact
        argument boundary of \(J_r\);
  \item every proof edge and boundary witness uniquely discharges one selected
        socket, with no missing or unused discharge;
  \item every proof node is proof-edge reachable from the root and the proof-edge
        graph is acyclic, with its topological order derived rather than stored;
  \item a supplying node is shared only when every incoming edge requires the
        identical formed judgment, including all arguments and aliasing equations;
        and
  \item every conclusion boundary, local witness, edge, socket, and boundary
        witness occurs exactly where projected, with no disconnected or unused
        position or incidence in \(d\).
\end{enumerate}

Proof-edge acyclicity is separate from the semantic formation-and-truth acyclicity
of Section~\ref{sec:audit}.  Neither supplies the other.

The no-unused conditions range over all positions and incidence of \(d\), not only
an \(h\)-selected subgraph.  Thus two passing selections use all the same proof
nodes, edges, sockets, boundaries, and witnesses, but may group or interpret them
through different actual-side images and project different premise judgments.
An alternative needing another position or incidence needs a different \(d\).

Conversely, one \((d,h)\) may pass for distinct supplied \(J_1,J_2\) when each
independently fits the generic catalogue and recursively required premises.  Each
validation derives its own premise meanings from the supplied judgment.  The
checker validates only the supplied \(h\); it neither enumerates other selections
nor requires their projections to agree.

\subsection{Assumptions and conditional derivation}

After a complete match, let
\[
  \operatorname{Assumptions}_G(J,d,h)
\]
be the exact finite family containing one
\(((k,r),J_r)\) for every selected premise socket discharged by a boundary
assertion.  The origin \((k,r)\) preserves multiplicity.  A boundary witness
contains its exact assertion subject, frame, version, frontier, and boundary
binding.  No ambient assumption registry contributes.

Define
\[
  \Derives_G(\Gamma,J,d,h)
  \Longleftrightarrow
  \CompleteMatch_G(J,d,h)
  \land
  \Gamma=\operatorname{Assumptions}_G(J,d,h).
\]
\(\Gamma\) is neither stored in \(d\) nor independently selected.  It exposes the
exact boundary-discharge projection of supplied \((J,d,h)\).  Another selection,
or another supplied judgment independently validated by the same \((d,h)\), may
project a different \(\Gamma\).

The dependency order is therefore
\[
\begin{array}{rcl}
\operatorname{LocalMatch}_F
 &\longrightarrow& \text{complete raw }m,\\
\CandidateAt_G
 &\longrightarrow& \text{arguments, prospective premises, and }\chi,\\
\MatchAt_G
 &\longrightarrow& \text{resolved formation, lookup, and checker success},\\
\CompleteMatch_G
 &\longrightarrow& \text{actual edge/boundary discharge incidence},\\
\operatorname{Assumptions}_G
 &\longrightarrow& \text{the exact }\Gamma.
\end{array}
\]

Define closed validity and mathematical proof existence only now:
\[
\begin{split}
 \ValidDerivation_G(J,d,h)
 &\Longleftrightarrow \Derives_G(\varnothing,J,d,h),\\
 \Proves_G(J)
 &\Longleftrightarrow
 \exists\text{ finite }d,h:\ValidDerivation_G(J,d,h).
\end{split}
\]
\(\Proves_G\) is a mathematical existence statement, not semantic truth and not a
search procedure.  Reference validation checks supplied \((d,h)\); failure to
supply one proves no negation.

Replacing a boundary assertion at one socket by a valid proof edge is the
structural change that closes that assumption.  No flag or premise class does so.
Exact reuse uses an identical-judgment socket: an edge closes it, a boundary
assertion leaves it conditional, and a self-edge or cycle fails acyclicity.

\subsection{Soundness, termination, and proof locality}

\begin{theorem}[Conditional derivation soundness]\label{thm:derivation-soundness}
For every generic form, a passing instance has a true conclusion whenever all
formed judgments at its open premises are true.  Consequently,
\[
\begin{split}
 \Derives_G(\Gamma,J,d,h)
 \land
 \bigwedge_{((k,r),J_r)\in\Gamma}
       (J_r\text{ is true in }G)
 &\Longrightarrow J\text{ is true in }G,\\
 \ValidDerivation_G(J,d,h)
 &\Longrightarrow J\text{ is true in }G,\\
 \Proves_G(J)
 &\Longrightarrow J\text{ is true in }G.
\end{split}
\]
\end{theorem}

\begin{proof}
Each primitive row checks its displayed semantic equation directly.  Trace,
enumeration, composite, negative, pattern, reuse, and parametric forms have the
soundness clauses stated in Section~\ref{sec:proof-forms}; the parametric case also
uses truth of its exact admissibility premise.  Induct over the topological order of
the proof edges selected by \(h\).  Edge-discharged premises obtain truth from
earlier nodes.  Boundary-discharged premises remain precisely the hypotheses in
\(\Gamma\).  The root then has truth \(J\).
\end{proof}

\begin{theorem}[Finite derivation checking]\label{thm:derivation-termination}
Checking supplied finite \((G,J,d,h)\) terminates.  It performs no search for a
missing node or alternative selection, invokes no semantic truth oracle, and
enumerates no unbounded schema domain.
\end{theorem}

\begin{theorem}[Layer-2 locality]\label{thm:proof-locality}
If two Factor models agree on every raw formation input, complete incident star,
four-component declared bundle, fixed endpoint, and positive or absent restricted
binding result for every raw match at every selected image, then they agree on
\(\CompleteMatch_G\), \(\operatorname{Assumptions}_G\), and \(\Derives_G\),
including false results.  This statement contains no retained image, realization
site, access trace, or implementation strategy.
\end{theorem}

Permitted renaming of owners wholly internal to \(d\) transports \(h\), each local
binding, premise edge, boundary witness, and assumption origin while fixing exact
semantic endpoints in \(G\), and preserves \(\Derives_G\).

\subsection{Derivation examples}

One finite \(d\) may contain two symmetric complete local-evidence substructures
and discharge positions arranged so that two distinct selections each account for
every position and incidence.  One selection can project a boundary assertion at
the first socket and an edge at the second; the other can reverse those
interpretations.  They yield different conditional \(\Gamma\) values without a
stored mode.  If either reading needed an extra witness position, it would require
a different \(d\).

The same actual selection may also fit two supplied judgments whose
\(\operatorname{ArgBoundary}_G\) values coincide, provided each semantic target
independently satisfies a passing generic-form instance and all recursively
projected premises.  The supplied judgment, not \(h\), asks the semantic question.

\section{Structural realization and normative replay}\label{sec:replay}

\subsection{Closed retained format and structural projection}

This layer uses the same fixed model \(G\).  Model formation is already a
metalevel precondition; it is not replay evidence.

There is one closed proposal-authored realization formula.  It contains no
semantic clause, proof rule, premise meaning, assumption truth, or proof checker.
Its root criterion is
\[
\begin{split}
 \operatorname{EnvelopeShape}_G(p,z;C,F,Q,L)\Longleftrightarrow
 \exists t\,[{}&\Rec_G(p;t,C,F,Q,L)\land\RefPos_G(t;z)\\
 &{}\land\operatorname{ConclusionImage}_G(C,z)
 \land\operatorname{FragmentFamily}_G(F)\\
 &{}\land\operatorname{PathGraph}_G(Q;p,z,F)
 \land\operatorname{SelectorFamily}_G(L,F)].
\end{split}
\]
The four payload predicates are finite substrate equations:

\begin{itemize}[leftmargin=2em]
  \item \(\operatorname{ConclusionImage}_G(C,z)\) says that \(C\) owns exactly
        one complete unlabelled conclusion-boundary image and its one external
        target reference is \(z\);
  \item \(\operatorname{FragmentFamily}_G(F)\) says that \(F\) is an exact
        \(\Fam\) of reference positions to distinct retained fragment roots, each
        root owning one complete fragment in the layer-2 proof grammar;
  \item \(\operatorname{PathGraph}_G(Q;p,z,F)\) says that \(Q\) is an exact
        \(\Fam\) of two-reference \(\Rec\) edge records, every edge is ordinary
        ownership or \(\operatorname{End}\) incidence in \(G\), every target and
        fragment root is reachable from \(p\), and every listed edge occurs in at
        least one such path; and
  \item \(\operatorname{SelectorFamily}_G(L,F)\) says that \(L\) is an exact
        \(\Fam\) of two-reference \(\Rec\) link records whose endpoints lie in
        the fragment-owned proof image.
\end{itemize}

Every family member and endpoint is quantified by these equations, and
\(\Rec_G(p;t,C,F,Q,L)\) forbids a sixth field.  Thus the envelope has the following
exact meaning-bearing components:

\begin{itemize}[leftmargin=2em]
  \item one target binding to exact actual root position \(z\), together with the
        retained image of the common unlabelled conclusion boundary;
  \item finite retained proof fragments and their structurally exposed roots;
  \item one shared finite path subgraph from \(p\) to every exposed target and
        fragment root; and
  \item complete/no-extra ordinary owned selector-link positions whose endpoints
        are retained proof-image positions.
\end{itemize}

Each path step is ordinary ownership incidence or an actual reference with its
exact endpoint.  Paths may share prefixes.  A path is not duplicated for internal
positions of a fragment.  Every exposed root is reachable, every path edge
contributes to a reachability, and every retained position projected as abstract
proof structure occurs under exactly one fragment root by ordinary ownership and
incidence.  Target-envelope, path, and selector-link positions are each accounted
for by their own equations.  Merely traversing such scaffolding does not make it a
proof-image position.  Missing, overlapping, arbitrary, or unused structure
rejects the format.

Let \(I_{G,F}\) be the union of the exact owned closures of the fragment roots,
restricted to positions satisfying one of the closed layer-2 proof-record
equations.  Formation requires those closures to be pairwise disjoint and every
position under a fragment root to be accepted by exactly one proof-record
equation.  Define the structural projection relation
\[
  \ProjectsAt_G(c,p,z;\bar d,\bar h,e)
\]
to hold exactly when \(c\) is admitted, \(p\in\operatorname{Body}_G(c)\), there
are unique \(C,F,Q,L\) satisfying \(\operatorname{EnvelopeShape}_G(p,z;C,F,Q,L)\),
and:

\begin{enumerate}[leftmargin=2.2em]
  \item \(e\) is a bijection
\[
  e:\operatorname{Pos}(\bar d)\longrightarrow\operatorname{Pos}(G)
\]
        with range exactly \(I_{G,F}\);
  \item transporting the restrictions of
        \(\operatorname{Own}_G,\operatorname{End}_G,
        \operatorname{Next}_G,\operatorname{Bind}_G\) on \(I_{G,F}\) backward
        through \(e\) gives exactly the primitive tables of \(\bar d\), including
        every common conclusion boundary and no scaffold position; and
  \item \(\bar h\) is exactly
        \[
          \{(e^{-1}(x),e^{-1}(y)):
             \text{one and only one link record in }L
             \text{ has endpoints }x,y\}.
        \]
\end{enumerate}

This is a restriction, inverse image, and finite equality of primitive tables.  It
does not invoke a decoder or a judgment of sameness of meaning.  Incidences are
transported as relations rather than mapped as separate objects.  Envelope, path,
or selector positions alone create no member of \(\operatorname{range}(e)\).

No canonical representative or owner order is chosen.  Exact occurrence
identities, \(c,p,z\), graph-backed external endpoints, and every boundary-crossing
owner remain fixed and outside \(\operatorname{Pos}(\bar d)\).  The projection checks
only complete/no-extra realization structure.  It does not resolve a local proof
match, validate a proof edge, or interpret a semantic judgment.

\begin{theorem}[Structural-projection determinacy]\label{thm:projection}
If two outputs project from the same \((G,c,p,z)\), there is one unique permitted
structural isomorphism
\(\phi:\bar d_1\cong\bar d_2\) transporting \(\bar h_1\) to \(\bar h_2\) and
satisfying
\[
  e_2(\phi(q))=e_1(q)
  \qquad(q\in\operatorname{Pos}(\bar d_1)).
\]
The origin maps make \(\phi\) unique even when the abstract proof has internal
automorphisms.
\end{theorem}

\subsection{The realization correspondence}

Let \(\iota:d\cong\bar d\) be a permitted incidence-preserving structural
isomorphism that fixes every external endpoint in \(G\).  The argument \(b\) is the
finite metalevel correspondence
\[
  b=\{(q,e(\iota(q))):q\in\operatorname{Pos}(d)\}.
\]
Its domain is exactly \(\operatorname{Pos}(d)\).  It preserves and reflects
ownership, ordered endpoint incidence, multiplicity, identity, aliasing, sharing,
conclusion boundaries, and open endpoints.  Graph-backed endpoints in \(G\) are
fixed by \(\iota\) rather than included in \(b\); binder-owned hypothetical
positions belong to \(\operatorname{Pos}(d)\) and are included.

The storage representation does not contain \(b\).  It retains only the ordinary
structures quantified by \(\operatorname{EnvelopeShape}\) and \(\ProjectsAt\).
Changing abstract representatives transports
\(\iota,e,b,\bar h\) together and changes no realization fact.

For every \((k,q)\in h\), the selector subgraph contains exactly one ordinary
owned link whose endpoints are \(b(k)\) and \(b(q)\), and contains no other
selector link.  Equivalently, \(\iota\) transports \(h\) exactly to \(\bar h\).
Selector links contain no catalogue pattern position, local binding, form
identity, or copy of a proof pattern.

\subsection{Pure structural realization}

Define
\[
  \RealizedAt(G;J,d,h,z,c,p,b).
\]
\(z\) is one actual graph-backed owner/local root position, not a set or metalevel
target.  The occurrence \(c\) is admitted and
\(p\in\operatorname{Body}_G(c)\).  The envelope at
\(p\) satisfies one \(\operatorname{EnvelopeShape}\), and its proof image is
projected by \(\ProjectsAt_G\).  The relation holds exactly when:

\begin{enumerate}[leftmargin=2.2em]
  \item \(\Formed_G(J)\), \(d\) has one structural root \(k_0\), and
        \(C_{k_0}\) is structurally equal to
        \(\operatorname{ArgBoundary}_G(J)\);
  \item the envelope binds the projected image of that boundary exactly to \(z\),
        preserving every supplied argument, graph endpoint, identity, aliasing,
        sharing, order, and incidence equation, including mapped finite
        binder-owned hypothetical structure;
  \item there exist \(\bar d,\bar h,e,\iota\) with
        \(\ProjectsAt_G(c,p,z;\bar d,\bar h,e)\),
        \(\iota:d\cong\bar d\) exact on \(\operatorname{Pos}(d)\), every external
        endpoint fixed, and \(b=\{(q,e(\iota(q)))\}\) with no missing or extra pair;
  \item \(\iota\) transports \(h\) exactly to \(\bar h\); and
  \item fragments may be distributed among \(c\) and referenced occurrences, but
        every conclusion boundary, proof edge, socket, local witness, and boundary
        witness is reconstructed exactly and every fragment root is reached by the
        shared path.
\end{enumerate}

\(\RealizedAt\) is pure layout and correspondence.  It performs no resolved local
matching, complete derivation check, proof-edge acyclicity check, soundness check,
or assumption-truth check.  It may hold for a structurally realized invalid
\((d,h)\).  A retained fragment can participate in several realization envelopes.
For a fixed envelope, different supplied \((d,h,b)\) witnesses differ only by the
permitted origin-preserving isomorphism.

\subsection{Containing occurrences}

For an actual stored position \(s\), define
\[
  \operatorname{ContainingOccurrence}_G(s)
\]
by following the exact finite ownership chain to the nearest admitted occurrence
whose body contains that chain.  If none exists or uniqueness fails, the
projection is undefined.  When \(s\) itself is an admitted occurrence root, it is
its own containing occurrence.  No occurrence-kind tag participates.

\(b\) maps ordinary positions only.  Proof-edge and link incidence is represented
by ordinary owned positions and preserved as a relation.  Thus every
\(s\in\operatorname{range}(b)\) is a legitimate argument of this projection.

\subsection{The sole retained-validity relation}

Define
\[
\begin{split}
 &\Replay(G;\Gamma,J,d,h,z,c,p,b)\\
 &\quad\Longleftrightarrow\
 \Derives_G(\Gamma,J,d,h)\\
 &\qquad{}\land\RealizedAt(G;J,d,h,z,c,p,b)\\
 &\qquad{}\land
   \operatorname{ContainingOccurrence}_G(z)<_{\mathrm{anc},G}c\\
 &\qquad{}\land
   \forall s\in\operatorname{range}(b):\
   \left[
     \operatorname{ContainingOccurrence}_G(s)=c
     \ \lor\
     \operatorname{ContainingOccurrence}_G(s)<_{\mathrm{anc},G}c
   \right].
\end{split}
\]
This is the complete conjunct list.  The path selects fragment roots; ordinary
ownership reconstructs their internal images.  The final conjuncts validate only
the containing occurrences and ancestry of selected positions.  Supplying proofs
are not recursively replayed because every proof node and edge already occurs in
finite \(d\), and \(h\) already fixes the reading.  Every mapped boundary assertion
is retained, but replay does not assert its truth.

An exact tuple satisfying this relation is called a \emph{certified realization}
in prose.  It is closed when \(\Gamma=\varnothing\) and conditional otherwise.
There is no certificate sort or second certified relation.  Existence claims
quantify directly over \(\Replay\).  Failure to supply an exact tuple establishes
no negative fact and grants no permission to search for one.

\begin{theorem}[Replay consequences and termination]\label{thm:replay}
A closed replay implies
\(\ValidDerivation_G(J,d,h)\), \(\Proves_G(J)\), and truth of \(J\) in \(G\).
A conditional replay preserves its exact assumption projection and does not turn a
boundary assertion into truth.  Reference replay terminates on finite supplied
\((G,d,h,b)\).
\end{theorem}

\subsection{Finite locality without a normative footprint}

Replay has no dependency-footprint argument or output.  For each supplied tuple,
its truth or falsehood is nevertheless determined by finite structural inputs:

\begin{itemize}[leftmargin=2em]
  \item judgment and prospective-premise formation, including presentation or
        schema boundaries and finite semantic audits;
  \item every raw match at a selected image, its four declared components, and
        every resolved comparison;
  \item every explicit envelope/projection alternative, realization structure,
        the supplied \(b\) boundary,
        containing-owner chains, admission causes, frontiers, and ancestry; and
  \item complete incident stars, fixed endpoints, and positive or absent restricted
        binding results used by those checks.
\end{itemize}

This finite locality does not evaluate semantic truth merely because a semantic
target occurs in a formation audit.  It never enumerates substitutions quantified
by \(\UniversalAdmissible_G\).  The nonnormative appendix gives a constructive
sufficient-footprint proof.

\begin{theorem}[Two-sided replay locality]\label{thm:replay-locality}
If Factor models \(G,G'\) agree on one sufficient finite structural footprint of
the kind above, including all fixed endpoints and restricted binding results, then
the same supplied tuple has the same \(\Replay\) truth value in both models.  Both
models independently satisfy metalevel model formation.
\end{theorem}

An implementation may use the appendix footprint, a conservative finite superset,
the whole finite model, or full recomputation.  Extra positions chosen for
invalidation do not become semantic dependencies.

\subsection{Retained external assumptions}

For a passing derivation and realization, let the unique complete finite value
\(\operatorname{BoundaryImage}_{G,z,c,p,b}(k,r)\) contain the exact stored images
of the premise socket and its finite boundary witness, the relevant envelope/path
positions, and the projection-origin correspondence from which those images follow.
Let
\[
  \operatorname{RetainedOrigin}_{b}(k,r)=(b(k),b(r)).
\]
Define the public projection
\[
\begin{split}
 &\operatorname{ExternalReq}
       (G;\Gamma,J,d,h,z,c,p,b)\\
 &\quad=
 \left\{
 \left(
  \operatorname{RetainedOrigin}_{b}(k,r),
  \bigl(J_r,\operatorname{BoundaryImage}_{G,z,c,p,b}(k,r)\bigr)
 \right):
 ((k,r),J_r)\in\Gamma
 \right\}.
\end{split}
\]
There is no tag or tagged union.  Injectivity of \(b\) keeps equal-looking
assertions at different sockets distinct.  Edge-discharged premises have no
external entry, and no recursive union is needed because the complete reading
already contains every supplying node.

\begin{theorem}[Complete transport of external requirements]%
\label{thm:external-transport}
Let \(\rho\) be a permitted renaming internal to \(d\).  Transport \(d,h,(k,r)\)
and bound positions of each \(J_r\), fix every exact graph endpoint, and transport
the realization witnesses so \(b'(\rho(q))=b(q)\).  Then retained boundary images
remain the same exact graph positions while abstract coordinates transport by
\(\rho\) and the induced structural-projection isomorphism.  The public origin stays
\((b(k),b(r))\), and the complete \(\operatorname{ExternalReq}\) value is equal
under that transport.
\end{theorem}

\subsection{Direct retained use and usable adoption}

A raw authority claim and a replayed realization become connected only through a
supplied actual consumer position.  All four arguments below are supplied graph
positions.  Define
\[
  \operatorname{DirectUseAt}_G(c,p,u,x)
\]
by the exact equation
\[
\begin{split}
 \operatorname{DirectUseAt}_G(c,p,u,x)\Longleftrightarrow{}&
 (\exists f\,\operatorname{AdmittedAt}_G(c,f))
       \land p,u\in\operatorname{Body}_G(c)\\
 &{}\land\exists a,b\,[\Rec_G(u;a,b)
                    \land\RefPos_G(a;p)\land\RefPos_G(b;x)].
\end{split}
\]
Thus the complete structure rooted at \(u\) owns exactly two ordinary reference
positions in encoded order with endpoint tuple \((p,x)\).  There is no missing or
extra owned position, reference, or endpoint.  The tuple is actual ordered
incidence, not a stored site/object role or direction field.

The consumer at \(u\) refers to \(p\) through that tuple, but it is outside the realization
envelope and is not quantified by \(\operatorname{EnvelopeShape}\) or
\(\ProjectsAt\).  The relation is finite and
renaming invariant while fixing its supplied graph positions.  It is functional:
if the same \(c,u\) passes with \((p,x)\) and \((p',x')\), then \(p=p'\) and
\(x=x'\).  Two \(\Rec\) witnesses with disagreeing endpoint projections make the
relation false.

\begin{theorem}[Direct-use exactness]\label{thm:direct-use}
\(\operatorname{DirectUseAt}_G\) is finite and invariant under every permitted
internal renaming that fixes its supplied graph positions.  If both
\(\operatorname{DirectUseAt}_G(c,p,u,x)\) and
\(\operatorname{DirectUseAt}_G(c,p',u,x')\) hold, then \(p=p'\) and \(x=x'\).
Two complete \(\Rec\) witnesses for \(u\) that disagree on either endpoint make
the relation false rather than selecting one by priority.
\end{theorem}

Define
\[
\begin{split}
 &\operatorname{UsableAdoption}_{A,G}
       (\xi,\mathit{claim},z,d,h,c,p,u,b)\\
 &\quad\Longleftrightarrow
 \exists J,\Gamma\text{ finite}:\
   \Formed_G(J)\\
 &\qquad{}\land
   \operatorname{ClaimCore}_{A,G}(\xi,\mathit{claim})=(J,z)\\
 &\qquad{}\land
   \operatorname{Adopts}_{A,G}(\xi,\mathit{claim},J,z)\\
 &\qquad{}\land
   \Replay(G;\Gamma,J,d,h,z,c,p,b)\\
 &\qquad{}\land
   \operatorname{DirectUseAt}_G(c,p,u,\mathit{claim}).
\end{split}
\]
\(\Gamma\) is unique because replay fixes it as
\(\operatorname{Assumptions}_G(J,d,h)\).  The use is unconditional exactly when
it is empty.  Proving the same \((J,z)\) without a direct link to the exact adopted
claim is insufficient.

When a replayed transition, release, or bridge result is consumed, its caller
likewise supplies an actual \(u\), checks the raw direct semantic relation, and
applies \(\operatorname{DirectUseAt}_G\) to the exact used object.  Mapped support
occurrences precede the carrying realization.  Both ordered bridge endpoints
precede its realization but need not be ancestors of one another.  Ancestry
establishes retention only and never reconstructs a semantic edge or direction.

\subsection{Realization examples}

Suppose \(J\) has one graph-backed argument and one binder-owned hypothetical
argument.  A realization fixes the first argument's exact endpoint in \(G\), maps
every owned position of the second through \(b\), and can retain two proof
fragments whose paths share one prefix from \(p\).  The shared prefix is scaffolding,
not duplicated proof structure.

If one retained premise socket changes from a boundary witness to an edge targeting
an independently valid supplying node, the new \(d,h\) has that origin removed
from \(\Gamma\) and from \(\operatorname{ExternalReq}\).  No assumption flag
changes.  Likewise, moving direct-use structure can change usable adoption while
the exact raw claim adoption remains unchanged.

\section{Implementation boundary and reference equivalence}\label{sec:implementation}

Encoding tags, hashes, compiled recognizers, caches, indexes, dispatch metadata,
database layout, and host-language types have no semantic, derivation,
realization, or replay force.  An implementation may validate or incrementally
maintain the metalevel fact that an input graph is a Factor model before invoking
the proposal's relations.  It may cache valid derivations, compile the fixed local
checkers, index retained citations and external requirements, reconstruct transient
\(b\), memoize realization comparisons, or invalidate on any conservative set.

The model component of the input at this boundary is the literal tuple
\[
(U_G,\operatorname{Own}_G,\operatorname{End}_G,
 \operatorname{Next}_G,\operatorname{Bind}_G).
\]
Every supplied finite argument
\(A\) contributes its literal carrier and restrictions of those same primitive
tables, including all boundary-crossing endpoints and restricted binding results;
supplied finite relations such as \(h,b,\theta\) contribute their literal pair
sets.  Let \(\mathcal I=(G;\vec A;\vec M)\) denote that whole mathematical input.
It is not bytes and is not the output of a proposal-owned decoder.  First write
\[
\begin{aligned}
 G\equiv_{\Sigma_F}G'\quad\Longleftrightarrow\quad{}
 &U_G=U_{G'}
 \land\operatorname{Own}_G=\operatorname{Own}_{G'}\\
 &{}\land\operatorname{End}_G=\operatorname{End}_{G'}
 \land\operatorname{Next}_G=\operatorname{Next}_{G'}\\
 &{}\land\operatorname{Bind}_G=\operatorname{Bind}_{G'}.
\end{aligned}
\]
and define
\[
 \mathcal I\equiv_{\Sigma_F}\mathcal I'
 \Longleftrightarrow
 G\equiv_{\Sigma_F}G'
 \land|\vec A|=|\vec A'|
 \land\bigwedge_{i=1}^{|\vec A|}
       \operatorname{Tables}(A_i)=\operatorname{Tables}(A_i')
 \land|\vec M|=|\vec M'|
 \land\bigwedge_{j=1}^{|\vec M|}M_j=M_j'.
\]
This is literal extensional equality of every finite input table and relation, not
merely equality of roots or host values.  A renamed input instead supplies one
explicit bijection preserving and reflecting every primitive tuple of \(G\) and
every \(A_i\), and transports each \(M_j\); it does not ask whether two encodings
have the same meaning.

\subsection{Canonical encoding-profile contract}

Factor does not choose a byte or host representation, but an artifact claiming to
be a \emph{canonical Factor encoding profile} must meet a fixed metalevel
contract.  A profile \(P\) publishes an exact specification and version, an
external carrier \(W_P\) that includes malformed as well as well-formed finite
words, and a declared nonempty supported domain
\(D_P\) of complete finite inputs \(\mathcal I\).  A constrained profile states
its domain constraints; it may not silently treat an unsupported input as a
different supported one.  The exact profile, not merely a profile name, is agreed
before parsing.

The profile supplies total single-valued maps
\[
\begin{aligned}
 \operatorname{Parse}_P &: W_P\longrightarrow
                 D_P\mathbin{\sqcup}\{\mathsf{reject}\},\\
 \operatorname{Encode}_P &: D_P\longrightarrow W_P .
\end{aligned}
\]
The published grammar must determine both maps by terminating effective
procedures.  Total parsing means that every external word terminates with either
one complete table-bundle result or the one explicit result
\(\mathsf{reject}\); truncation, an unknown field, an out-of-domain atom, or
ambiguity never leaves interpretation undefined.  Total encoding terminates on
all of \(D_P\).  For every
\(\mathcal I,\mathcal I'\in D_P\) and \(w\in W_P\), the profile laws are
\[
\begin{aligned}
 \operatorname{Parse}_P(\operatorname{Encode}_P(\mathcal I))
     &\equiv_{\Sigma_F}\mathcal I,\\
 \operatorname{Parse}_P(w)=\mathcal I
     &\Longrightarrow \operatorname{Encode}_P(\mathcal I)=w,\\
 \operatorname{Encode}_P(\mathcal I)
     =\operatorname{Encode}_P(\mathcal I')
     &\Longleftrightarrow
       \mathcal I\equiv_{\Sigma_F}\mathcal I'.
\end{aligned}
\]
The first is exact round trip, the second rejects noncanonical alternatives, and
the third is injectivity together with extensional stability.  A profile that
intentionally accepts several surface forms instead supplies a total
normalization into such a canonical profile; those surface forms are not
themselves the canonical encoding.

The specification must account for every component used in
\(\mathcal I\equiv_{\Sigma_F}\): carrier identities; the arity and order of
\(\vec A,\vec M\); every
\(\operatorname{Own}\), \(\operatorname{End}\), \(\operatorname{Next}\), and
\(\operatorname{Bind}\) tuple; the restricted primitive tables and every
boundary-crossing endpoint of each supplied finite argument; and every pair,
including the empty pair set, of each supplied relation.  It must preserve
distinct incidence positions, repeated endpoints, aliasing, sharing, encoded
order, and absent restricted bindings.  It may not consult an ambient registry,
locale, mutable default, host class, or omitted-field convention, silently
collapse duplicate rows, or synthesize a missing table.  Rejection forms no
Factor query and establishes no negative Factor fact.

For accepted words under one exact profile, define
\[
 w\equiv_P w'
 \Longleftrightarrow
 \exists\mathcal I,\mathcal I'\,
 [\operatorname{Parse}_P(w)=\mathcal I
  \land\operatorname{Parse}_P(w')=\mathcal I'
  \land\mathcal I\equiv_{\Sigma_F}\mathcal I'].
\]
This is the wire-level meaning of the phrase \emph{the same input}.  It is
defined by the published maps and literal table equality, not by similarity of
names or intention.  For two different profiles \(P,Q\), byte equality has no
force: their accepted results must first be compared by
\(\equiv_{\Sigma_F}\), or transported by one explicit primitive-table
isomorphism.

\begin{lemma}[Canonical-profile same-input criterion]\label{lem:profile}
For \(w,w'\) accepted by one canonical profile \(P\),
\(w\equiv_P w'\) if and only if \(w=w'\).  Two adapters proved conforming to that
same \(P\) therefore return extensionally equal inputs, or both return
\(\mathsf{reject}\), on the same external word.
\end{lemma}

\begin{proof}
If the parsed inputs are extensionally equal, injectivity makes their canonical
encodings equal, and the second profile law identifies those encodings with
\(w,w'\).  The converse follows from single-valued parsing.  The adapter statement
then follows by equality with the same total \(\operatorname{Parse}_P\).
\end{proof}

For example, a profile may restrict its complete input carrier to
\(\{0,\ldots,n-1\}\).  One length-delimited word records \(n\), then the four
lexicographically ordered primitive relation tables, then a length-delimited block
for every supplied argument and relation.  Zero lengths encode empty
\(\operatorname{Bind}\) and pair tables.  Its parser rejects unsorted or duplicate
rows, out-of-range atoms, wrong lengths, omitted blocks, and trailing material;
its encoder emits the unique ordered form.  The example assigns no meaning to a
tag: it simply gives one auditable representation of the literal tables for its
declared domain.

A concrete adapter \(D\) is \emph{\(P\)-conforming} exactly when
\(D(w)=\operatorname{Parse}_P(w)\) for every \(w\in W_P\), including rejection
cases.  Testing or proving this equality is external implementation work.  Neither
the profile maps nor their conformance is a Factor relation.  Two incompatible
parse maps are two different profiles even if their documents reuse one identifier;
their adapters cannot both satisfy the same \(P\)-conformance claim.

An implementation is \emph{reference-conforming} exactly when, for every Factor
input \(\mathcal I\) presented as those tables and every exact argument tuple,
each exposed result equals the authoritative mathematical relation on that tuple.  This is an
external extensional criterion, not a Factor judgment or an internally generated
certificate.  No \(\operatorname{Conforms}\) relation is defined; determining
whether arbitrary implementation code meets the criterion is outside Factor.  The
compared results are:

\begin{itemize}[leftmargin=2em]
  \item semantic formation and truth for every relation it exposes;
  \item complete raw local matches and their four-component declared inputs;
  \item resolved \(\MatchAt_G\), complete derivation, and assumption projections;
  \item \(\operatorname{EnvelopeShape}\), \(\ProjectsAt\), and realization
        results; and
  \item \(\Replay\) truth and every derived semantic output.
\end{itemize}

Dependency-set identity, minimality, access order, and cache-invalidation strategy
are not conformance outputs.  No cache or index supplies a missing match,
derivation, assumption, correspondence, direct semantic edge, or ancestry edge.

Only here may \emph{build} abbreviate \(\Apply\) followed by an external storage
action.  It has no Factor relation of its own.  A physical encoding may materialize
a derived output body or copied input bodies for integrity checking, but those are
erasable caches: the reference body is derived from ancestry and the authoritative
cause.  Creation APIs, repository copying, byte encodings, capability tiers,
bootstrap installations, and index designs remain outside the proposal.

The following is a corollary of this definition and the structural-transport
theorems, not an internal certification theorem.

\begin{corollary}[Extensional implementation agreement]\label{thm:conformance}
Two reference-conforming implementations given
\(\mathcal I\equiv_{\Sigma_F}\mathcal I'\) return identical reference results.
If instead their complete inputs are related by one explicit primitive-table isomorphism,
their results are related by the induced structural transport.  Additional reads,
short-circuiting, physical layout, cached intermediates, and conservative
invalidation supersets cannot change either conclusion.
\end{corollary}

This corollary alone deliberately does not compare two adapters from a common byte
string, database row, host object, tag set, or serialization.  If two adapters are
separately proved \(P\)-conforming to the same exact canonical profile, the
Lemma~\ref{lem:profile} first supplies the same \(\mathcal I\), after which the corollary
applies.  If they produce different primitive tables, at least one adapter fails
that common profile claim or the adapters are using different profiles.  In the
latter case the proposal does not select an interpretation: the parsed table
bundles must satisfy \(\equiv_{\Sigma_F}\), or one explicit table isomorphism
must be supplied, before reference results can be compared.  These adapter and
profile premises are auditable external obligations, not Factor facts silently
assumed by the corollary.

A positive finite fact and a complete negative fact remain different.  For
example, one exact mismatching endpoint proves two supplied endpoints unequal.
To prove that no member of an explicit finite domain matches a candidate, the
complete-negative form supplies the whole duplicate-free domain and one positive
mismatch for every member.  Failure of an index or search proves neither claim.

\boxtext{\textbf{Authority of definitions.}  The metalevel model conditions and
Sections 1--7 give the complete proposal.
Later appendices prove consequences and give nonnormative bookkeeping; they add no
Factor relation or competing validity condition.}

\appendix

\section{Nonnormative proof of replay locality}\label{app:locality}

\textbf{Status.}  This appendix is nonnormative.  It constructs one sufficient
finite set for the proof of Theorem~\ref{thm:replay-locality}.  The set has no
public Factor name, is not a replay input or output, and need not be minimal,
unique, canonical, retained, or returned by an implementation.

For one supplied tuple, begin with a finite worklist and finite sets of graph
positions, key values, and restricted binding results.

\begin{enumerate}[leftmargin=2.2em]
  \item Add the complete inputs deciding formation of the root judgment and every
        structurally reached prospective premise: application boundaries,
        presentation or schema boundaries, raw finite meaning carriers, dependency
        edges, and the acyclicity audit.  Do not add semantic truth values or
        positions used only to evaluate a truth clause.
  \item At every selected actual image, add every complete raw match of every fixed
        catalogue clause, all four declared components, complete incident stars,
        resolved comparisons, fixed checker operands, and each independently
        formed child edge.  Preserve candidates that fail later; do not stop at the
        first failure.
  \item Add every \(\operatorname{EnvelopeShape}\) and \(\ProjectsAt\)
        alternative, envelope boundary, fragment root,
        shared-path and selector-link position, supplied \(b\) boundary, and exact
        target comparison needed to preserve success or failure of realization.
  \item Add every containing-owner chain, admission cause, authoritative body
        dependency, frontier, and recursively reached parent boundary required to
        preserve both positive and negative ancestry checks.
  \item Whenever one of those fixed equations quantifies over a complete incident
        star, add the whole star.  Add every exact endpoint.  For every reached key
        add its positive or absent restricted binding result and, when positive,
        its endpoint.  Continue until the finite worklist is empty.
\end{enumerate}

Finiteness follows because \(G,d,h,b\), every raw pattern, and every fixed
formation carrier are finite; each closure step adds positions from those finite
objects.  The construction never expands a semantic closure, enumerates the
possibly infinite substitutions in \(\UniversalAdmissible_G\), or evaluates a
truth clause solely because its definition target is audited.

Agreement of two Factor models on the resulting structural boundary preserves
the functional presentation predicates and their acyclicity, every raw candidate and
declared bundle, every resolved success or failure, and every prospective premise.
In particular, agreement on a restricted binding with no pair preserves that
failed lookup as absence; no first-failure path can turn it into a positive result.
Induction over the supplied proof-edge order preserves
\(\CompleteMatch_G\), \(\operatorname{Assumptions}_G\), and \(\Derives_G\).
The envelope/projection alternatives and origin maps then agree up to their unique
origin-preserving isomorphism.  Owner-chain and frontier agreement preserves
containment and ancestry in both directions.  The complete replay conjunction
therefore has the same truth value.

Processing all alternatives simultaneously proves independence from checker order,
short-circuiting, and a failure in a sibling or parent candidate.  A source failure
cannot suppress an independently formed child.  An implementation may use a larger
set or the entire model; those extra positions do not become semantic dependencies.
Exact retained citations and \(\operatorname{ExternalReq}\) are structural outputs,
whereas this worklist is only proof bookkeeping.

\section{Proofs of structural metatheorems}\label{app:proofs}

\subsection{Model, ownership, and meaning}

Finite acyclicity of the frontier-parent relation makes strict ancestry well
founded: every nonempty descending chain would repeat a frontier and form a cycle.
This proves the recursion premise in Theorem~\ref{thm:body}.  Existence follows by
well-founded recursion; uniqueness follows because each base cause carries one body
and each derived cause has one complete \(V\) and deterministic \(\Apply\).
Renaming transports every owner, endpoint, and map equation bijectively, proving
Theorem~\ref{thm:kernel-renaming}.

\(\operatorname{Pos}(A)\) is functional because transitive ownership in finite
structure has an exact least closure.  External endpoints are never owned and
binder-owned descendants always are, so permitted renaming preserves and reflects
the boundary.  The formation equations then determine one
\(\operatorname{AppBoundary}_G\).  The incidence-derived seeds and least closure
determine one \(\operatorname{ArgBoundary}_G\); the exclusion of target-only
incidence proves that equal argument projections do not merge semantic targets.

For Theorem~\ref{thm:meaning}, Lemma~\ref{lem:presentation-separation} separates
the three presentation patterns even when the extensional tuple family and unary
closure seed have the same shape.  In the structural-view case,
Lemma~\ref{lem:clause-separation} proves pairwise disjointness of all eight clause
rows, including the singleton-family projection case; exact boundary exhaustion
forbids a missing row, hybrid, or residual marker.  Functionality then follows
from the exact \(\Rec\) spines, reflected payload grammars, and agreement condition
on two \(\ExactMatch\) witnesses.  Root equality transports complete boundaries
and therefore transports each fixed outgoing edge.  Identity, inverse, and
composition show that the stated graph-backed-or-hypothetical root comparison is an
equivalence.  The same transport preserves each incidence-defined
\(\operatorname{DefOcc}_G(a)\) and its unique outgoing dependency set.  The root
carrier is bounded by finite \(G\) plus finitely supplied hypothetical roots; the
occurrence carrier is bounded by those roots and finitely many proposal clauses.
Their union is finite, so reachability and cycle detection are decidable.  Acyclic
dependency recursion supplies one formation and truth meaning.

\subsection{Semantic structural relations}

\(\SameStructure_G\) is an equivalence by identity, inverse, and composition of
its exact bijections.  The proof of Theorem~\ref{thm:congruence} transports each
finite \(\Apply\) operand through jointly compatible node, position, and boundary
partial bijections.  They preserve resolved requirements and injective matches,
and extending the maps by the fresh pairing yields the output correspondence.
Root-erasing equality alone lacks the external-boundary equations and therefore
does not support this transport.

Least closure supplies the trace characterization of reachability by induction on
finite trace length in one direction and on membership derivations in the other.
Factorization glue is deterministic because the complete interface has one inverse
rewrite per cut endpoint.  Generated-family equality closure stays inside
\(\mathcal D_{G,X}\) by its explicit guards.  Induction on its least closure proves
the usual seed-and-preservation theorem used by necessity.

Property projection is deterministic because the formed application boundary
fixes the exact subject position, forced alias class, remaining arguments, and
simultaneous enclosing equations.  The exact target incidence creates one
\(\operatorname{DefOcc}_G(a)\); merging occurrences for other targets would change
that incidence and is forbidden by the audit equality.

\subsection{Finite domains and admissibility}

Identity, inverse, and composition of complete root-boundary bijections prove that
\(\operatorname{MemberIso}_G\) is an equivalence.  Every outside graph position is
fixed in both directions, so symmetry cannot swap external endpoints.  Finite
enumeration of root-preserving bijections proves decidability.  Duplicate-free
payloads then make the matching representative unique.

The six least-boundary rules for \(\operatorname{SchemaBoundary}_G\) are monotone
over finite \(G\) and terminate at one least fixed point.  Traversal of each finite
image similarly determines \(\operatorname{FullImageBoundary}_G\).
Exact-domain/range and incidence reflection make \(\beta_\theta\) functional.
Formation of instantiated boundaries is a finite total check.  Model agreement on
these complete operands preserves every conjunct of
\(\operatorname{PermBind}_G\), including its failure when an instantiation is
unformed.  The unique \(\operatorname{SchemaShape}\) projection,
duplicate-free guard formation, and finite
cycle detection make admissibility formation functional, decidable, and
renaming-invariant.  The audit includes all fixed hypothesis and conclusion
targets, so admissibility cannot approve itself through a meaning cycle.  These
facts prove Theorem~\ref{thm:schema}.

\subsection{Catalogue, matching, and derivations}

The DF tables enumerate every primitive operator and every
permitted carrier variant.  Each raw pattern is finite and each input profile is a
least closure over finite structure.  The remaining generic forms have the same
property.  This proves termination and exhaustiveness in
Theorem~\ref{thm:catalogue}; a clause absent from the tables has no direct
applicability equation.  In the parametric form, its finite substitution and
boundary checks establish one permitted \(\theta\); truth of the exact
admissibility premise and every instantiated open hypothesis then gives the
instantiated conclusion by the universal clause.  Thus the form's substitution
theorem invokes no installed rule or conclusion oracle.

For a fixed completed raw match, the row fixes all seeds and closure rules, so
\(\operatorname{DeclaredInputs}\) is functional.  Equal images transport these
seeds and equations and therefore have coherent projections.  Resolved checks add
conjuncts without changing the raw binding or bundle, proving
Theorem~\ref{thm:local-determinacy}.  Exact agreement of complete stars and
restricted binding results preserves each conjunct in both directions, proving
Theorem~\ref{thm:match-locality}.

Every complete reading uses all positions of \(d\); acyclicity gives a finite
topological order.  Induction in that order proves
Theorem~\ref{thm:derivation-soundness}.  Local checker termination and finite graph
traversal prove Theorem~\ref{thm:derivation-termination}.  Applying
Theorem~\ref{thm:match-locality} independently at every selected image, including
all failed alternatives, preserves the complete discharge incidence and exact
assumption projection.  This proves Theorem~\ref{thm:proof-locality}.

\subsection{Projection, replay, external requirements, and authority}

The explicit proof-image restriction assigns every projected proof position one
retained origin.  Composing the two origin bijections yields the isomorphism in
Theorem~\ref{thm:projection}; equality at every retained origin makes it unique.
Composition with \(\iota\) gives an exact \(b\), and the selector equations
transport \(h\) without copying catalogue positions.

Realization is finite structural comparison.  Containment follows finite owner
chains, and ancestry follows finite frontier parents.  Conjoining these checks
with finite derivation checking proves replay termination.  Closed replay invokes
Theorem~\ref{thm:derivation-soundness}, proving
Theorem~\ref{thm:replay}.  Appendix~\ref{app:locality} proves the two-sided locality
claim without adding a semantic footprint.

Under a permitted internal renaming, \(b'(\rho(q))=b(q)\) fixes every public
retained origin.  The induced structural-projection isomorphism transports the
abstract coordinates while the graph images stay exact, proving
Theorem~\ref{thm:external-transport}.  The two-reference complete/no-extra
equation at \(u\) makes \(\operatorname{DirectUseAt}_G\) finite and functional:
two passing endpoint tuples must be the same ordered tuple.  Incidence-preserving
renaming transports that unique tuple while fixing its supplied graph positions,
which proves Theorem~\ref{thm:direct-use}.

Finally, raw claim-core and adoption projections read only their fixed frame,
claim, target, and raw support positions.  Realization and consumer positions lie
outside that projection, so changing only them preserves the raw result.  This
proves Theorem~\ref{thm:raw-authority}; usable adoption can still change because it
has the separate replay and direct-use conjuncts.

\subsection{Reference equivalence}

Every reference relation is a mathematical function or relation of the literal
primitive tables and supplied finite relations.  The definition of reference conformance
therefore gives identical outputs when
\(\mathcal I\equiv_{\Sigma_F}\mathcal I'\).  The proposal's
renaming theorems give the transported result for an explicit primitive-table
isomorphism.  A cache, type, index, or access trace is absent from every
authoritative argument and cannot justify another answer.  This yields
Corollary~\ref{thm:conformance}.  On its own it supplies no premise about
independently written encoding adapters.  Lemma~\ref{lem:profile} supplies that
premise only after each adapter has separately met the same published
\(P\)-conformance criterion; arbitrary or cross-profile adapters remain outside
the corollary.

\end{document}
